\chapter{Chapter09 — Matching Moment Combinatorics}
%%% generated by the blueprint pipeline from FABL.Chapter09.MatchingMomentCombinatorics
\section{Overview}

Declarations of the Lean module \texttt{FABL.Chapter09.MatchingMomentCombinatorics}.

\section{Declarations}

\begin{definition}[Perfect matching codes]
\lean{FABL.PerfectMatchingCode}
\statementsource{Analysis-of-Boolean-Functions-by-Ryan-ODonnell}{pdf-ffbe9c7b3306-p277-b004}
\label{FABL.PerfectMatchingCode}
\leanok
Define a family of types recursively by setting the family at $0$ equal to a singleton,
and the family at $m+1$ equal to the product of a set of $2m+1$ choices and the family
at $m$.
\end{definition}

\begin{definition}[Finite matching codes]
\lean{FABL.perfectMatchingCodeFintype}
\label{FABL.perfectMatchingCodeFintype}
\leanok
\uses{FABL.PerfectMatchingCode}
Define a type of codes for each natural number $m$ recursively: for $m=0$ it is a
singleton, and for $m+1$ it consists of a choice from $\{0,\ldots,2m\}$ together with a
code for $m$.
\end{definition}

\begin{definition}[Equality on matching codes]
\lean{FABL.perfectMatchingCodeDecidableEq}
\label{FABL.perfectMatchingCodeDecidableEq}
\leanok
\uses{FABL.PerfectMatchingCode}
For each natural number $m$, define a decision procedure for equality of matching codes:
at $m=0$ compare the singleton codes, and at $m+1$ compare the chosen index and the code
for $m$ componentwise.
\end{definition}

\begin{definition}[Decoding matching codes]
\lean{FABL.decodePerfectMatchingCode}
\label{FABL.decodePerfectMatchingCode}
\leanok
\uses{FABL.PerfectMatchingCode}
Given a type $\alpha$, a natural number $m$, a matching code of size $m$, and a list of
elements of $\alpha$, define a list of ordered pairs recursively. For $m=0$ the output
is empty; for $m+1$, if the input list has first element $x$ and the coded index selects
an element $y$ among the remaining entries, output $(x,y)$ followed by the recursively
decoded pairs after removing that selected entry; if no such entry exists, or the input
list is empty, the output is empty.
\end{definition}

\begin{definition}[Decoded matching]
\lean{FABL.decodedPerfectMatching}
\label{FABL.decodedPerfectMatching}
\leanok
\uses{FABL.PerfectMatchingCode, FABL.decodePerfectMatchingCode}
For a natural number $m$ and a matching code of size $m$, define its decoded matching by
decoding it against the list $0,1,\ldots,2m-1$.
\end{definition}

\begin{definition}[Partite matching code condition]
\lean{FABL.IsPartitePerfectMatchingCode}
\statementsource{Analysis-of-Boolean-Functions-by-Ryan-ODonnell}{pdf-ffbe9c7b3306-p277-b004}
\label{FABL.IsPartitePerfectMatchingCode}
\leanok
\uses{FABL.PerfectMatchingCode, FABL.decodePerfectMatchingCode}
For natural numbers $q,k$ and a perfect matching code of index $(qk)/2$, decode the code
on the list of vertices $0,1,\ldots,qk-1$. The code is partite when every decoded edge
has endpoints whose integer quotients upon division by $k$ are unequal.
\end{definition}

\begin{definition}[Partite perfect matchings]
\lean{FABL.PartitePerfectMatching}
\statementsource{Analysis-of-Boolean-Functions-by-Ryan-ODonnell}{pdf-ffbe9c7b3306-p277-b004}
\label{FABL.PartitePerfectMatching}
\leanok
\uses{FABL.IsPartitePerfectMatchingCode, FABL.PerfectMatchingCode}
For natural numbers $q,k$, define a partite perfect matching to be a perfect matching
code of index $(qk)/2$ satisfying the condition that every decoded edge joins vertices
with unequal quotients upon division by $k$.
\end{definition}

\begin{theorem}[Perfect matching code count]
\lean{FABL.natCard_perfectMatchingCode}
\label{FABL.natCard_perfectMatchingCode}
\leanok
\uses{FABL.PerfectMatchingCode}
\difficulty{4}
For every natural number $m$, the number of perfect matching codes with parameter $m$ is
$(2m-1)!!$.
\end{theorem}

\begin{theorem}[Partite matching code count bound]
\lean{FABL.natCard_partitePerfectMatching_le_code}
\label{FABL.natCard_partitePerfectMatching_le_code}
\leanok
\uses{FABL.PartitePerfectMatching, FABL.PerfectMatchingCode}
\difficulty{1}
For natural numbers $q$ and $k$, the number of partite perfect matching codes is at most
the number of perfect matching codes of length parameter $(qk)/2$.
\end{theorem}

\begin{theorem}[Four partite matchings on eight vertices]
\lean{FABL.card_partitePerfectMatching_four_two}
\label{FABL.card_partitePerfectMatching_four_two}
\leanok
\uses{FABL.PartitePerfectMatching}
\difficulty{1}
Among the perfect matchings on eight vertices divided into four parts of size two,
exactly sixty have every matched pair in different parts.
\end{theorem}

\begin{theorem}[Four partite matchings on four vertices]
\lean{FABL.card_partitePerfectMatching_four_one}
\label{FABL.card_partitePerfectMatching_four_one}
\leanok
\uses{FABL.PartitePerfectMatching}
\difficulty{1}
Among the perfect matchings on four vertices divided into four singleton parts, exactly
three have every matched pair in different parts.
\end{theorem}

\begin{theorem}[Multiplicity counterexample for partite matchings]
\lean{FABL.exercise9_37b_printed_multiplicity_counterexample}
\label{FABL.exercise9_37b_printed_multiplicity_counterexample}
\leanok
\uses{FABL.PartitePerfectMatching, FABL.card_partitePerfectMatching_four_one}
\difficulty{3}
The number of partite perfect matchings encoded for four parts of size one, divided by
$(1!)^4$, is not equal to $1$.
\end{theorem}

\begin{theorem}[Odd double factorial bound]
\lean{FABL.oddDoubleFactorial_le_factorial_sq}
\label{FABL.oddDoubleFactorial_le_factorial_sq}
\leanok
\difficulty{5}
For every natural number $m$, the double factorial $(2m-1)!!$ is at most $(m!)^2$.
\end{theorem}

\begin{theorem}[Affine moments by the binomial formula]
\lean{FABL.pmfExpectation_affine_pow}
\label{FABL.pmfExpectation_affine_pow}
\leanok
\uses{FABL.pmfExpectation}
\difficulty{3}
Let $\Omega$ be a finite set, let $\mathrm{law}$ be a probability mass function on
$\Omega$, let $X:\Omega\to\mathbb R$, let $a,c\in\mathbb R$, and let $m$ be a natural
number. Then $\mathbb
E_{\mathrm{law}}[(a+cX)^m]=\sum_{j=0}^{m}\binom{m}{j}a^{m-j}c^j\mathbb
E_{\mathrm{law}}[X^j]$.
\end{theorem}

\begin{theorem}[Odd-zero sum reduces to even terms]
\lean{FABL.sum_range_two_mul_add_one_eq_sum_even}
\label{FABL.sum_range_two_mul_add_one_eq_sum_even}
\leanok
\difficulty{4}
Let $R$ be a commutative additive monoid, let $F$ be a sequence in $R$, and let $r$ be a
nonnegative integer. Suppose $F(2j+1)=0$ for every natural number $j<r$. Then
\[\sum_{j=0}^{2r}F(j)=\sum_{j=0}^{r}F(2j).\]
\end{theorem}

\begin{theorem}[Even moment bound under moment conditions]
\lean{FABL.affine_evenMoment_le_of_momentConditions}
\label{FABL.affine_evenMoment_le_of_momentConditions}
\leanok
\uses{FABL.pmfExpectation, FABL.pmfExpectation_affine_pow, FABL.pmfExpectation_const, FABL.sum_range_two_mul_add_one_eq_sum_even}
\difficulty{6}
Let $\Omega$ be a finite nonempty set, let $\mu$ be a probability law on $\Omega$, and
let $X:\Omega\to\mathbb R$. Suppose $r$ is a positive integer, $\rho\ge0$, and
$\rho^2\le 1/(2r-1)$. Assume that $\mathbb E_\mu[X^2]=1$, that every odd moment of order
$2j-1$ vanishes for $1\le j\le r$, and that for each such $j$ the moment of order $2j$
is at most $\big((2r-1)^j\binom rj\big)/\binom{2r}{2j}$. Then, for all $a,b\in\mathbb
R$, $\mathbb E_\mu[(a+\rho bX)^{2r}]\le(a^2+b^2)^r$.
\end{theorem}

\begin{theorem}[Second norm of a centered affine function]
\lean{FABL.pmfLpNorm_two_affine_centered}
\label{FABL.pmfLpNorm_two_affine_centered}
\leanok
\uses{FABL.pmfExpectation, FABL.pmfExpectation_affine_pow, FABL.pmfExpectation_const, FABL.pmfLpNorm}
\difficulty{4}
Let $\Omega$ be a finite set, let $\mathrm{law}$ be a probability mass function on
$\Omega$, and let $X:\Omega\to\mathbb{R}$ have expectation $0$ and second moment $1$.
For all $a,b\in\mathbb{R}$, the $L^2$ norm of the function $x\mapsto a+bX(x)$ with
respect to $\mathrm{law}$ is $\sqrt{a^2+b^2}$.
\end{theorem}

\begin{definition}[Even hypercontractivity for a PMF]
\lean{FABL.IsPmfTwoEvenHypercontractive}
\label{FABL.IsPmfTwoEvenHypercontractive}
\leanok
\uses{FABL.pmfLpNorm}
Let $\Omega$ be a finite set, let $\mu$ be a probability mass function on $\Omega$, let
$r$ be a natural number, let $\rho\in\mathbb R$, and let $X:\Omega\to\mathbb R$. The
property holds when, for every $a,b\in\mathbb R$, the $L^{2r}$ norm under $\mu$ of the
function $x\mapsto a+\rho bX(x)$ is at most the $L^2$ norm under $\mu$ of $x\mapsto
a+bX(x)$.
\end{definition}

\begin{theorem}[Moment conditions imply hypercontractivity]
\lean{FABL.isPmfTwoEvenHypercontractive_of_momentConditions}
\label{FABL.isPmfTwoEvenHypercontractive_of_momentConditions}
\leanok
\uses{FABL.IsPmfTwoEvenHypercontractive, FABL.affine_evenMoment_le_of_momentConditions, FABL.pmfExpectation, FABL.pmfLpNorm, FABL.pmfLpNorm_nonneg, FABL.pmfLpNorm_rpow, FABL.pmfLpNorm_two_affine_centered}
\difficulty{5}
Let $\Omega$ be a finite set with probability mass function $\mu$, let
$X:\Omega\to\mathbb{R}$, let $r$ be a positive integer, and let $\rho\ge0$ satisfy
$\rho^2\le 1/(2r-1)$. Suppose $\mathbb{E}_{\mu}[X^2]=1$, $\mathbb{E}_{\mu}[X^{2j-1}]=0$
for every integer $j$ with $1\le j\le r$, and, for each such $j$,
$\mathbb{E}_{\mu}[X^{2j}]\le (2r-1)^j\binom{r}{j}/\binom{2r}{2j}$. Then for all
$a,b\in\mathbb{R}$, the $2r$-norm of $a+\rho bX$ is at most the $2$-norm of $a+bX$.
\end{theorem}

\begin{theorem}[Even powers of the probability norm]
\lean{FABL.pmfLpNorm_even_pow_eq_pmfExpectation_even}
\label{FABL.pmfLpNorm_even_pow_eq_pmfExpectation_even}
\leanok
\uses{FABL.pmfExpectation, FABL.pmfLpNorm, FABL.pmfLpNorm_nonneg, FABL.pmfLpNorm_rpow}
\difficulty{3}
Let $\Omega$ be a finite set, let $\pi$ be a probability distribution on $\Omega$, let
$r$ be a positive integer, and let $f:\Omega\to\mathbb R$. Define
$\|f\|_{p,\pi}=\left(\sum_{x\in\Omega}\pi(x)|f(x)|^p\right)^{1/p}$. Then
$\|f\|_{2r,\pi}^{2r}=\sum_{x\in\Omega}\pi(x)f(x)^{2r}$.
\end{theorem}

\begin{theorem}[Even affine moment from odd moments]
\lean{FABL.pmfExpectation_affine_evenMoment_of_oddMoments}
\label{FABL.pmfExpectation_affine_evenMoment_of_oddMoments}
\leanok
\uses{FABL.pmfExpectation, FABL.pmfExpectation_affine_pow, FABL.sum_range_two_mul_add_one_eq_sum_even}
\difficulty{4}
Let $\Omega$ be a finite set, let $\mathrm{law}$ be a probability mass function on
$\Omega$, let $X:\Omega\to\mathbb R$, let $r$ be a natural number, and let $c\in\mathbb
R$. Suppose that for every natural number $j$ with $1\le j\le r$, $\mathbb
E_{\mathrm{law}}[X^{2j-1}]=0$. Then $\mathbb
E_{\mathrm{law}}[(1+cX)^{2r}]=\sum_{j=0}^{r}\binom{2r}{2j}c^{2j}\mathbb
E_{\mathrm{law}}[X^{2j}]$.
\end{theorem}

\begin{theorem}[Necessary bound for hypercontractivity]
\lean{FABL.isPmfTwoEvenHypercontractive_necessary_sq}
\label{FABL.isPmfTwoEvenHypercontractive_necessary_sq}
\leanok
\uses{FABL.IsPmfTwoEvenHypercontractive, FABL.pmfExpectation, FABL.pmfExpectation_affine_evenMoment_of_oddMoments, FABL.pmfExpectation_const, FABL.pmfExpectation_nonneg, FABL.pmfLpNorm_even_pow_eq_pmfExpectation_even, FABL.pmfLpNorm_nonneg, FABL.pmfLpNorm_two_affine_centered}
\difficulty{6}
Let $\Omega$ be a finite set with probability mass function $\mu$, let
$X:\Omega\to\mathbb{R}$, let $r$ be a positive integer, and let $\rho\in\mathbb{R}$.
Suppose $\mathbb{E}_{\mu}[X^2]=1$, every odd moment $\mathbb{E}_{\mu}[X^{2j-1}]$
vanishes for $1\le j\le r$, and for all $a,b\in\mathbb{R}$ the $2r$-norm of $a+\rho bX$
is at most the $2$-norm of $a+bX$. Then $\rho^2\le 1/(2r-1)$.
\end{theorem}

\begin{theorem}[Hypercontractivity from moment conditions]
\lean{FABL.isPmfTwoEvenHypercontractive_iff_of_momentConditions}
\label{FABL.isPmfTwoEvenHypercontractive_iff_of_momentConditions}
\leanok
\uses{FABL.IsPmfTwoEvenHypercontractive, FABL.isPmfTwoEvenHypercontractive_necessary_sq, FABL.isPmfTwoEvenHypercontractive_of_momentConditions, FABL.pmfExpectation}
\difficulty{5}
Let $\Omega$ be a finite set, let $\pi$ be a probability mass function on $\Omega$, and
let $X:\Omega\to\mathbb R$. Let $r$ be a positive natural number and let $\rho\ge0$.
Suppose $\mathbb E_\pi[X^2]=1$, every odd moment $\mathbb E_\pi[X^{2j-1}]$ vanishes for
$1\le j\le r$, and every even moment satisfies $\mathbb E_\pi[X^{2j}]\le
\frac{(2r-1)^j\binom rj}{\binom{2r}{2j}}$ for $1\le j\le r$. Then the inequality
$\|a+\rho bX\|_{2r,\pi}\le\|a+bX\|_{2,\pi}$ holds for all real $a,b$ if and only if
$\rho\le 1/\sqrt{2r-1}$.
\end{theorem}

\begin{theorem}[Even moments of a uniform sign]
\lean{FABL.pmfExpectation_uniformSign_evenMoment}
\proofsource{Analysis-of-Boolean-Functions-by-Ryan-ODonnell}{
  pdf-ffbe9c7b3306-p257-b004,
  pdf-ffbe9c7b3306-p258-b003
}
\label{FABL.pmfExpectation_uniformSign_evenMoment}
\leanok
\uses{FABL.Sign, FABL.pmfExpectation, FABL.pmfExpectation_uniformPMF_eq_expect, FABL.signValue, FABL.uniformPMF}
\difficulty{3}
For every nonnegative integer $j$, the uniform expectation of $s^{2j}$ over
$s\in\{\pm1\}$ equals $1$.
\end{theorem}

\begin{theorem}[Odd moments of a uniform sign]
\lean{FABL.pmfExpectation_uniformSign_oddMoment}
\proofsource{Analysis-of-Boolean-Functions-by-Ryan-ODonnell}{pdf-ffbe9c7b3306-p257-b004}
\label{FABL.pmfExpectation_uniformSign_oddMoment}
\leanok
\uses{FABL.Sign, FABL.pmfExpectation, FABL.pmfExpectation_uniformPMF_eq_expect, FABL.signValue, FABL.uniformPMF}
\difficulty{3}
For every nonnegative integer $j$, the uniform expectation of $s^{2j+1}$ over
$s\in\{\pm1\}$ equals $0$.
\end{theorem}

\begin{theorem}[Uniform sign two even of rho le]
\lean{FABL.uniformSign_twoEven_of_rho_le}
\statementsource{Analysis-of-Boolean-Functions-by-Ryan-ODonnell}{
  pdf-ffbe9c7b3306-p258-b001,
  pdf-ffbe9c7b3306-p258-b003
}
\label{FABL.uniformSign_twoEven_of_rho_le}
\leanok
\uses{FABL.Sign, FABL.oneBit_twoQ, FABL.signValue, FABL.uniformLpNorm, FABL.uniformLpNorm_two_affineSign}
\difficulty{4}
% dataset-pipeline: stub -- prose pending build_informal_statements.py
Statement pending regeneration.
\end{theorem}

\begin{theorem}[Even moments of an affine uniform sign]
\lean{FABL.pmfExpectation_uniformSign_affine_evenMoment}
\proofsource{Analysis-of-Boolean-Functions-by-Ryan-ODonnell}{pdf-ffbe9c7b3306-p257-b004}
\label{FABL.pmfExpectation_uniformSign_affine_evenMoment}
\leanok
\uses{FABL.Sign, FABL.pmfExpectation, FABL.pmfExpectation_affine_pow, FABL.pmfExpectation_uniformSign_evenMoment, FABL.pmfExpectation_uniformSign_oddMoment, FABL.signValue, FABL.sum_range_two_mul_add_one_eq_sum_even, FABL.uniformPMF}
\difficulty{4}
For every nonnegative integer $r$ and real number $c$, the uniform expectation over
$s\in\{\pm1\}$ of $(1+c s)^{2r}$ is $\sum_{j=0}^{r}\binom{2r}{2j}c^{2j}$.
\end{theorem}

\begin{theorem}[Even moments of the uniform norm]
\lean{FABL.uniformLpNorm_even_pow_eq_expect_even}
\label{FABL.uniformLpNorm_even_pow_eq_expect_even}
\leanok
\uses{FABL.uniformLpNorm, FABL.uniformLpNorm_rpow}
\difficulty{4}
For a finite set $\Omega$, a positive natural number $r$, and a function
$f:\Omega\to\mathbb{R}$, the uniform $L^{2r}$ quantity of $f$, raised to the power $2r$,
equals the uniform average of $f(x)^{2r}$ over $\Omega$.
\end{theorem}

\begin{theorem}[Necessary noise bound for even exponents]
\lean{FABL.uniformSign_twoEven_necessary_sq}
\proofsource{Analysis-of-Boolean-Functions-by-Ryan-ODonnell}{pdf-ffbe9c7b3306-p257-b004}
\label{FABL.uniformSign_twoEven_necessary_sq}
\leanok
\uses{FABL.Sign, FABL.pmfExpectation_uniformPMF_eq_expect, FABL.pmfExpectation_uniformSign_affine_evenMoment, FABL.signValue, FABL.uniformLpNorm, FABL.uniformLpNorm_even_pow_eq_expect_even, FABL.uniformLpNorm_nonneg, FABL.uniformLpNorm_two_affineSign}
\difficulty{6}
Let $S=\{-1,1\}$ have the uniform probability distribution, let $X(s)=s$, let $r$ be a
positive integer, and let $\rho$ be real. Suppose that for every $a,b\in\mathbb{R}$, the
uniform $L^{2r}$ norm of $a+\rho bX$ is at most the uniform $L^2$ norm of $a+bX$. Then
$\rho^2\leq 1/(2r-1)$.
\end{theorem}

\begin{theorem}[Uniform sign two even iff]
\lean{FABL.uniformSign_twoEven_iff}
\statementsource{Analysis-of-Boolean-Functions-by-Ryan-ODonnell}{
  pdf-ffbe9c7b3306-p257-b004,
  pdf-ffbe9c7b3306-p258-b001
}
\label{FABL.uniformSign_twoEven_iff}
\leanok
\uses{FABL.Sign, FABL.signValue, FABL.uniformLpNorm, FABL.uniformSign_twoEven_necessary_sq, FABL.uniformSign_twoEven_of_rho_le}
\difficulty{4}
% dataset-pipeline: stub -- prose pending build_informal_statements.py
Statement pending regeneration.
\end{theorem}

\begin{definition}[Symmetric difference of a finite-set list]
\lean{FABL.listSymmDiff}
\statementsource{Analysis-of-Boolean-Functions-by-Ryan-ODonnell}{pdf-ffbe9c7b3306-p277-b004}
\label{FABL.listSymmDiff}
\leanok
For a list of finite subsets of a type, define its symmetric difference recursively: the
empty list has symmetric difference the empty set, and adjoining a first set takes its
symmetric difference with that set.
\end{definition}

\begin{definition}[Symmetric difference of a finite-set tuple]
\lean{FABL.tupleSymmDiff}
\statementsource{Analysis-of-Boolean-Functions-by-Ryan-ODonnell}{pdf-ffbe9c7b3306-p277-b004}
\label{FABL.tupleSymmDiff}
\leanok
\uses{FABL.listSymmDiff}
Given a natural number $q$ and a family of finite subsets indexed by
$\\{0,\\ldots,q-1\\}$, define their symmetric difference to be the symmetric difference
of the list formed from that family in index order.
\end{definition}

\begin{theorem}[Membership and even occurrence]
\lean{FABL.not_mem_listSymmDiff_iff_even_countP}
\label{FABL.not_mem_listSymmDiff_iff_even_countP}
\leanok
\uses{FABL.listSymmDiff}
\difficulty{4}
Let $\alpha$ be a type with decidable equality, let $x\in\alpha$, and let $\mathcal S$
be a finite list of finite subsets of $\alpha$. Then $x$ is absent from the symmetric
difference of the sets in $\mathcal S$ if and only if the number of sets in the list
that contain $x$ is even.
\end{theorem}

\begin{theorem}[Even membership count from empty symmetric difference]
\lean{FABL.even_countP_of_tupleSymmDiff_eq_empty}
\label{FABL.even_countP_of_tupleSymmDiff_eq_empty}
\leanok
\uses{FABL.listSymmDiff, FABL.not_mem_listSymmDiff_iff_even_countP, FABL.tupleSymmDiff}
\difficulty{3}
Let $\alpha$ be a type with decidable equality, let $q$ be a nonnegative integer, and
let $S_0,\ldots,S_{q-1}$ be finite subsets of $\alpha$. If their iterated symmetric
difference is empty, then for every $x\in\alpha$, the number of these sets that contain
$x$ is even.
\end{theorem}

\begin{theorem}[Indicator sum as a predicate count]
\lean{FABL.sum_ite_eq_countP_ofFn}
\label{FABL.sum_ite_eq_countP_ofFn}
\leanok
\difficulty{6}
Let $q$ be a nonnegative integer and let $P$ be a predicate on $\\{0,\\ldots,q-1\\}$.
The sum, over these indices, of $1$ when $P(i)$ holds and $0$ otherwise equals the
number of entries satisfying $P$ in the ordered list $0,1,\\ldots,q-1$.
\end{theorem}

\begin{theorem}[Counting elements satisfying a predicate]
\lean{FABL.natCard_subtype_eq_countP_ofFn}
\label{FABL.natCard_subtype_eq_countP_ofFn}
\leanok
\uses{FABL.sum_ite_eq_countP_ofFn}
\difficulty{3}
Let $q$ be a natural number and let $p$ be a decidable predicate on
$\{0,1,\ldots,q-1\}$. The number of indices $i$ in this set for which $p(i)$ holds
equals the number of true instances of $p$ in the list $0,1,\ldots,q-1$.
\end{theorem}

\begin{theorem}[Product of cube monomials]
\lean{FABL.list_prod_monomial}
\label{FABL.list_prod_monomial}
\leanok
\uses{FABL.listSymmDiff, FABL.monomial, FABL.monomial_mul_monomial}
\difficulty{4}
For any finite list $S_1,\ldots,S_k$ of subsets of $\{1,\ldots,n\}$ and any sign vector
$x\in\{\pm1\}^n$, the product of the associated cube monomials equals the monomial
indexed by the symmetric difference of all the sets: $\prod_{j=1}^k\prod_{i\in
S_j}x_i=\prod_{i\in S_1\triangle\cdots\triangle S_k}x_i$.
\end{theorem}

\begin{theorem}[Product of monomials under symmetric difference]
\lean{FABL.prod_monomial_eq_tupleSymmDiff}
\proofsource{Analysis-of-Boolean-Functions-by-Ryan-ODonnell}{
  pdf-ffbe9c7b3306-p024-b005,
  pdf-ffbe9c7b3306-p024-b006
}
\label{FABL.prod_monomial_eq_tupleSymmDiff}
\leanok
\uses{FABL.list_prod_monomial, FABL.monomial, FABL.tupleSymmDiff}
\difficulty{3}
Let $q,n$ be natural numbers, let $S_0,\ldots,S_{q-1}$ be subsets of the $n$
coordinates, and let $x$ be a sign vector on those coordinates. The product of the
monomials indexed by the sets $S_j$, evaluated at $x$, equals the monomial indexed by
their iterated symmetric difference, evaluated at $x$.
\end{theorem}

\begin{theorem}[Moment expansion by Fourier coefficients]
\lean{FABL.expect_pow_eq_sum_fourierCoeff_of_tupleSymmDiff_empty}
\proofsource{Analysis-of-Boolean-Functions-by-Ryan-ODonnell}{
  pdf-ffbe9c7b3306-p277-b004,
  pdf-ffbe9c7b3306-p024-b007,
  pdf-ffbe9c7b3306-p025-b001
}
\label{FABL.expect_pow_eq_sum_fourierCoeff_of_tupleSymmDiff_empty}
\leanok
\uses{FABL.expect_monomial, FABL.fourierCoeff, FABL.fourier_expansion, FABL.monomial, FABL.prod_monomial_eq_tupleSymmDiff, FABL.tupleSymmDiff}
\difficulty{6}
Let $n,q$ be natural numbers and let $f$ be a real-valued function on sign vectors of
length $n$. The uniform expectation of $f(x)^q$ equals the sum, over all $q$-tuples of
subsets of the coordinates whose symmetric difference is empty, of the product of the
corresponding Fourier coefficients of $f$.
\end{theorem}

\begin{theorem}[Moment bound by Fourier coefficients]
\lean{FABL.expect_pow_le_sum_abs_fourierCoeff_of_tupleSymmDiff_empty}
\proofsource{Analysis-of-Boolean-Functions-by-Ryan-ODonnell}{pdf-ffbe9c7b3306-p277-b004}
\label{FABL.expect_pow_le_sum_abs_fourierCoeff_of_tupleSymmDiff_empty}
\leanok
\uses{FABL.expect_pow_eq_sum_fourierCoeff_of_tupleSymmDiff_empty, FABL.fourierCoeff, FABL.tupleSymmDiff}
\difficulty{1}
Let $n,q$ be natural numbers and let $f$ be a real-valued function on sign vectors of
length $n$. The uniform expectation of $f(x)^q$ is at most the sum, over all $q$-tuples
of subsets of the coordinates whose symmetric difference is empty, of the absolute value
of the product of the corresponding Fourier coefficients of $f$.
\end{theorem}

\begin{definition}[Vertices of a complete multipartite graph]
\lean{FABL.BookPartiteVertex}
\statementsource{Analysis-of-Boolean-Functions-by-Ryan-ODonnell}{pdf-ffbe9c7b3306-p277-b004}
\label{FABL.BookPartiteVertex}
\leanok
For natural numbers $q$ and $k$, let the vertex set be the disjoint union of $q$ parts,
each containing $k$ elements; a vertex records both its part and its position within
that part.
\end{definition}

\begin{definition}[Complete multipartite graph]
\lean{FABL.bookPartiteGraph}
\statementsource{Analysis-of-Boolean-Functions-by-Ryan-ODonnell}{pdf-ffbe9c7b3306-p277-b004}
\label{FABL.bookPartiteGraph}
\leanok
\uses{FABL.BookPartiteVertex}
For natural numbers $q$ and $k$, define the graph on $q$ parts of $k$ vertices each in
which two vertices are adjacent exactly when they belong to different parts.
\end{definition}

\begin{definition}[Perfect matchings of a multipartite graph]
\lean{FABL.GraphPartitePerfectMatching}
\statementsource{Analysis-of-Boolean-Functions-by-Ryan-ODonnell}{pdf-ffbe9c7b3306-p277-b004}
\label{FABL.GraphPartitePerfectMatching}
\leanok
\uses{FABL.bookPartiteGraph}
For natural numbers $q$ and $k$, a graph-partite perfect matching is a subgraph of the
complete multipartite graph with $q$ parts of size $k$ that is a perfect matching.
\end{definition}

\begin{definition}[Labelings constant on matching edges]
\lean{FABL.MatchingCompatibleLabeling}
\statementsource{Analysis-of-Boolean-Functions-by-Ryan-ODonnell}{pdf-ffbe9c7b3306-p277-b004}
\label{FABL.MatchingCompatibleLabeling}
\leanok
\uses{FABL.BookPartiteVertex, FABL.bookPartiteGraph}
Let $q$ and $k$ be natural numbers, let $\\beta$ be a type, let each vertex of the
complete multipartite graph on $q$ parts of size $k$ receive a label in $\\beta$, and
let $M$ be a subgraph of that graph. The labeling is compatible with $M$ when the
endpoints of every edge of $M$ have equal labels.
\end{definition}

\begin{theorem}[Compatible perfect matching in a multipartite graph]
\lean{FABL.exists_graphPartitePerfectMatching_compatible}
\label{FABL.exists_graphPartitePerfectMatching_compatible}
\leanok
\uses{FABL.BookPartiteVertex, FABL.GraphPartitePerfectMatching, FABL.MatchingCompatibleLabeling, FABL.bookPartiteGraph}
\difficulty{5}
Let $q,k$ be natural numbers, and label the vertices of the complete $q$-partite graph
with $k$ vertices in each part by elements of a type $\beta$. Suppose the labeling is
injective on each part, and every label has an even number of vertices carrying it. Then
the graph has a perfect matching in which the two endpoints of every matched edge have
the same label.
\end{theorem}

\begin{theorem}[Weight bound from injective lifts]
\lean{FABL.sum_le_inv_mul_sum_of_injective_lifts}
\label{FABL.sum_le_inv_mul_sum_of_injective_lifts}
\leanok
\difficulty{4}
Let $A$ and $B$ be finite sets, let $m$ be a positive integer, and let $w:A\to\mathbb R$
and $v:B\to\mathbb R$ satisfy $v(b)\ge0$ for every $b\in B$. Suppose an injective map
$L:A\times[m]\to B$ satisfies $v(L(a,i))=w(a)$ for every $(a,i)\in A\times[m]$. Then
$\sum_{a\in A}w(a)\le \frac1m\sum_{b\in B}v(b)$.
\end{theorem}

\begin{theorem}[Partite matching double factorial bound]
\lean{FABL.natCard_partitePerfectMatching_le_doubleFactorial}
\label{FABL.natCard_partitePerfectMatching_le_doubleFactorial}
\leanok
\uses{FABL.PartitePerfectMatching, FABL.PerfectMatchingCode, FABL.natCard_partitePerfectMatching_le_code, FABL.natCard_perfectMatchingCode}
\difficulty{3}
For natural numbers $r$ and $k$, the number of partite perfect matching codes with
parameters $2r$ and $k$ is at most $(2rk-1)!!$, the double factorial of $2rk-1$.
\end{theorem}

\begin{theorem}[Partite matching factorial bound]
\lean{FABL.natCard_partitePerfectMatching_le_factorial_bound}
\label{FABL.natCard_partitePerfectMatching_le_factorial_bound}
\leanok
\uses{FABL.PartitePerfectMatching, FABL.natCard_partitePerfectMatching_le_doubleFactorial, FABL.oddDoubleFactorial_le_factorial_sq}
\difficulty{4}
Let $r$ and $k$ be natural numbers with $r>0$. The number of partite perfect matching
codes with parameters $2r$ and $k$, regarded as a real number, is at most
$\left(\frac{2r-1}{r}\right)^{rk}((rk)!)^2$.
\end{theorem}

\begin{theorem}[Homogeneous uniform lp norm le sqrt two mul]
\lean{FABL.homogeneous_uniformLpNorm_le_sqrt_two_mul}
\statementsource{Analysis-of-Boolean-Functions-by-Ryan-ODonnell}{
  pdf-ffbe9c7b3306-p277-b004,
  pdf-ffbe9c7b3306-p278-b001
}
\label{FABL.homogeneous_uniformLpNorm_le_sqrt_two_mul}
\leanok
\uses{FABL.fourierDegree, FABL.lowDegree_uniformLpNorm_le, FABL.uniformLpNorm, FABL.uniformLpNorm_nonneg}
\difficulty{3}
% dataset-pipeline: stub -- prose pending build_informal_statements.py
Statement pending regeneration.
\end{theorem}

\begin{theorem}[Multinomial coefficient bound]
\lean{FABL.multinomial_const_le_pow}
\label{FABL.multinomial_const_le_pow}
\leanok
\difficulty{4}
For natural numbers $r$ and $k$, the multinomial coefficient with $r$ parts, each equal
to $k$, satisfies $\binom{rk}{k,\ldots,k} \le r^{rk}$.
\end{theorem}

\begin{theorem}[Factorial product bound]
\lean{FABL.factorial_mul_le_factorial_pow_mul_pow}
\label{FABL.factorial_mul_le_factorial_pow_mul_pow}
\leanok
\uses{FABL.multinomial_const_le_pow}
\difficulty{4}
For all natural numbers $r$ and $k$, $(rk)!\le (k!)^r r^{rk}$.
\end{theorem}

\begin{definition}[Matching improvement constant]
\lean{FABL.matchingImprovementConstant}
\label{FABL.matchingImprovementConstant}
\leanok
For natural numbers $r$ and $k$, define the matching improvement constant by
\[\left(\frac{(rk)!}{(k!)^r r^{rk}}\right)^{1/(2r)},\] where the exponentiation is
real-valued.
\end{definition}

\begin{theorem}[Improvement constant is at most one]
\lean{FABL.matchingImprovementConstant_le_one}
\label{FABL.matchingImprovementConstant_le_one}
\leanok
\uses{FABL.factorial_mul_le_factorial_pow_mul_pow, FABL.matchingImprovementConstant}
\difficulty{4}
For natural numbers $r$ and $k$ with $r>0$, the improvement constant is at most $1$.
\end{theorem}

\begin{definition}[Disjoint block choices]
\lean{FABL.DisjointBlockChoices}
\label{FABL.DisjointBlockChoices}
\leanok
Let $\alpha$ be a type with decidable equality, let $R$ be a finite subset of $\alpha$,
and let $d,t\in\mathbb{N}$. Define the type of ordered sequences of $t$ pairwise
disjoint $d$-element subsets of $R$: for $t=0$ it has one element, and for $t+1$ an
element consists of a $d$-element subset of $R$ together with such a sequence chosen
from the elements of $R$ outside that subset.
\end{definition}

\begin{definition}[List of chosen blocks]
\lean{FABL.DisjointBlockChoices.blocks}
\label{FABL.DisjointBlockChoices.blocks}
\leanok
\uses{FABL.DisjointBlockChoices}
Given a type $\alpha$ with decidable equality, a finite set $R\subseteq\alpha$, a
natural number $d$, a natural number $t$, and a choice of $t$ pairwise disjoint
$d$-element subsets of $R$, define its block list to be the ordered list of those
subsets.
\end{definition}

\begin{theorem}[Length of the chosen block list]
\lean{FABL.DisjointBlockChoices.blocks_length}
\label{FABL.DisjointBlockChoices.blocks_length}
\leanok
\uses{FABL.DisjointBlockChoices, FABL.DisjointBlockChoices.blocks}
\difficulty{4}
For a type $\alpha$ with decidable equality, a finite set $R \subseteq \alpha$, natural
numbers $d$ and $t$, and any recursively defined choice of $t$ pairwise disjoint
$d$-element subsets of $R$, the associated list of blocks has length exactly $t$.
\end{theorem}

\begin{theorem}[Size of each chosen block]
\lean{FABL.DisjointBlockChoices.blocks_card}
\label{FABL.DisjointBlockChoices.blocks_card}
\leanok
\uses{FABL.DisjointBlockChoices, FABL.DisjointBlockChoices.blocks}
\difficulty{4}
Let $\alpha$ be a type with decidable equality, let $R$ be a finite subset of $\alpha$,
and let $d,t$ be natural numbers. For any valid sequence of $t$ disjoint choices of
$d$-element subsets, every block in the resulting list has cardinality $d$.
\end{theorem}

\begin{theorem}[Chosen blocks lie in the remaining set]
\lean{FABL.DisjointBlockChoices.blocks_subset}
\label{FABL.DisjointBlockChoices.blocks_subset}
\leanok
\uses{FABL.DisjointBlockChoices, FABL.DisjointBlockChoices.blocks}
\difficulty{4}
For a finite set of elements, a finite subset called the remaining set, and natural
numbers $d$ and $t$, consider a choice of $t$ pairwise disjoint blocks, each of size
$d$, drawn successively from the elements still remaining. Every block in the resulting
list is a subset of the original remaining set.
\end{theorem}

\begin{theorem}[Pairwise disjointness of chosen blocks]
\lean{FABL.DisjointBlockChoices.blocks_pairwiseDisjoint}
\label{FABL.DisjointBlockChoices.blocks_pairwiseDisjoint}
\leanok
\uses{FABL.DisjointBlockChoices, FABL.DisjointBlockChoices.blocks, FABL.DisjointBlockChoices.blocks_subset}
\difficulty{5}
Let $\alpha$ be a type with decidable equality, let $R$ be a finite subset of $\alpha$,
and let $d,t$ be natural numbers. The blocks in every valid sequence of $t$ disjoint
choices of $d$-element subsets are pairwise disjoint.
\end{theorem}

\begin{theorem}[Block choices are determined by their blocks]
\lean{FABL.DisjointBlockChoices.ext_of_blocks_eq}
\label{FABL.DisjointBlockChoices.ext_of_blocks_eq}
\leanok
\uses{FABL.DisjointBlockChoices, FABL.DisjointBlockChoices.blocks}
\difficulty{4}
Let a finite set of elements with decidable equality, a finite remaining set, and
natural numbers $d$ and $t$ be given. If two choices of $t$ pairwise disjoint blocks of
size $d$, drawn successively from the elements still remaining, have the same ordered
list of blocks, then the choices are equal.
\end{theorem}

\begin{definition}[Six disjoint block choices]
\lean{FABL.SixDisjointBlockChoice}
\label{FABL.SixDisjointBlockChoice}
\leanok
\uses{FABL.DisjointBlockChoices}
For natural numbers $n$ and $d$, define the type of six disjoint block choices to be the
type of choices of six ordered, pairwise disjoint subsets of $\{0,\ldots,n-1\}$, each of
size $d$.
\end{definition}

\begin{definition}[Finite structure on block choices]
\lean{FABL.disjointBlockChoicesFintype}
\label{FABL.disjointBlockChoicesFintype}
\leanok
\uses{FABL.DisjointBlockChoices}
For a type $\alpha$ with decidable equality, a finite set $R\subseteq\alpha$, and
natural numbers $d$ and $t$, define a finite type structure on the type of choices of
$t$ ordered, pairwise disjoint $d$-element subsets of $R$.
\end{definition}

\begin{definition}[Indexed block of a six-block choice]
\lean{FABL.SixDisjointBlockChoice.block}
\label{FABL.SixDisjointBlockChoice.block}
\leanok
\uses{FABL.DisjointBlockChoices.blocks, FABL.DisjointBlockChoices.blocks_length, FABL.SixDisjointBlockChoice}
For natural numbers $n$ and $d$, a six-block choice of $d$-element subsets of $in n$,
and an index $irom in 6$, define its $i$th block to be the entry at position $i$ in
the choice’s list of six blocks.
\end{definition}

\begin{theorem}[Cardinality of a chosen block]
\lean{FABL.SixDisjointBlockChoice.block_card}
\label{FABL.SixDisjointBlockChoice.block_card}
\leanok
\uses{FABL.DisjointBlockChoices.blocks_card, FABL.SixDisjointBlockChoice, FABL.SixDisjointBlockChoice.block}
\difficulty{3}
Let $n,d$ be natural numbers, and choose six pairwise disjoint $d$-element subsets of
the $n$-element set. Each of the six chosen blocks has cardinality $d$.
\end{theorem}

\begin{theorem}[Pairwise disjointness of chosen blocks]
\lean{FABL.SixDisjointBlockChoice.block_pairwiseDisjoint}
\label{FABL.SixDisjointBlockChoice.block_pairwiseDisjoint}
\leanok
\uses{FABL.DisjointBlockChoices.blocks_pairwiseDisjoint, FABL.SixDisjointBlockChoice, FABL.SixDisjointBlockChoice.block}
\difficulty{4}
Let $n,d$ be natural numbers, and choose six pairwise disjoint $d$-element blocks from
the $n$-element set. Any two blocks with distinct indices are disjoint.
\end{theorem}

\begin{theorem}[Cardinality of fixed-size subsets]
\lean{FABL.natCard_powersetCardSubtype}
\label{FABL.natCard_powersetCardSubtype}
\leanok
\difficulty{3}
Let $\alpha$ be a type and let $R$ be a finite subset of $\alpha$. For every natural
number $d$, the number of $d$-element subsets of $R$ is $\binom{|R|}{d}$.
\end{theorem}

\begin{theorem}[Count of disjoint block choices]
\lean{FABL.natCard_disjointBlockChoices}
\label{FABL.natCard_disjointBlockChoices}
\leanok
\uses{FABL.DisjointBlockChoices}
\difficulty{5}
Let $\alpha$ be a type with decidable equality, let $R$ be a finite subset of $\alpha$,
and let $d,t\in\mathbb N$. The number of ways to choose, in order, $t$ pairwise disjoint
subsets of $R$, each of size $d$, is $\displaystyle\prod_{j=0}^{t-1}\binom{|R|-jd}{d}$.
\end{theorem}

\begin{theorem}[Six disjoint block choice count]
\lean{FABL.natCard_sixDisjointBlockChoice}
\label{FABL.natCard_sixDisjointBlockChoice}
\leanok
\uses{FABL.SixDisjointBlockChoice, FABL.natCard_disjointBlockChoices}
\difficulty{1}
For natural numbers $n$ and $d$, the number of ways to choose six ordered, pairwise
disjoint blocks of size $d$ from an $n$-element set is $\prod_{j=0}^{5}\binom{n-jd}{d}$.
\end{theorem}

\begin{theorem}[Factorial identity for successive choices]
\lean{FABL.prod_choose_sub_mul_factorials}
\label{FABL.prod_choose_sub_mul_factorials}
\leanok
\difficulty{5}
For natural numbers $n,d,t$ satisfying $td\le n$, the identity
$\left(\prod_{j=0}^{t-1}\binom{n-jd}{d}\right)(d!)^t(n-td)!=n!$ holds.
\end{theorem}

\begin{definition}[Six-block multinomial coefficient]
\lean{FABL.sixBlockMultinomial}
\label{FABL.sixBlockMultinomial}
\leanok
For natural numbers $n$ and $d$, define the six-block multinomial coefficient as the
multinomial coefficient for seven categories whose sizes are $d$ in each of the first
six categories and $n-6d$ in the seventh category.
\end{definition}

\begin{theorem}[Six disjoint blocks counted by a multinomial coefficient]
\lean{FABL.natCard_sixDisjointBlockChoice_eq_multinomial}
\label{FABL.natCard_sixDisjointBlockChoice_eq_multinomial}
\leanok
\uses{FABL.SixDisjointBlockChoice, FABL.natCard_sixDisjointBlockChoice, FABL.prod_choose_sub_mul_factorials, FABL.sixBlockMultinomial}
\difficulty{5}
Let $n,d\in\mathbb{N}$ satisfy $6d\le n$. The number of ways to choose six ordered,
pairwise disjoint subsets of size $d$ from an $n$-element set is the multinomial
coefficient with six parts of size $d$ and one remaining part of size $n-6d$.
\end{theorem}

\begin{definition}[Edges of the complete four-vertex graph]
\lean{FABL.k4EdgeEnds}
\label{FABL.k4EdgeEnds}
\leanok
Define a map from the six indices $0,1,2,3,4,5$ to ordered pairs of vertices in
$\{0,1,2,3\}$ by listing, in order, $(0,1),(0,2),(0,3),(1,2),(1,3),(2,3)$.
\end{definition}

\begin{definition}[Edges incident to a vertex of the complete graph on four vertices]
\lean{FABL.k4IncidentEdges}
\label{FABL.k4IncidentEdges}
\leanok
\uses{FABL.k4EdgeEnds}
For each vertex of the complete graph on four vertices, let its incident-edge set
consist of all six edge indices whose corresponding edge has that vertex as an endpoint.
\end{definition}

\begin{theorem}[Degree of a vertex in the complete graph]
\lean{FABL.card_k4IncidentEdges}
\label{FABL.card_k4IncidentEdges}
\leanok
\uses{FABL.k4IncidentEdges}
\difficulty{4}
In the complete graph on four vertices, each vertex is incident to exactly three edges.
\end{theorem}

\begin{definition}[Incident support of a block choice]
\lean{FABL.SixDisjointBlockChoice.incidentSupport}
\label{FABL.SixDisjointBlockChoice.incidentSupport}
\leanok
\uses{FABL.SixDisjointBlockChoice, FABL.SixDisjointBlockChoice.block, FABL.k4IncidentEdges}
Given six successive choices of pairwise disjoint subsets of size $d$ from an
$n$-element set, and a vertex $v$ of the complete graph on four vertices, define its
incident support to be the union of the chosen subsets indexed by edges incident to $v$.
\end{definition}

\begin{theorem}[Cardinality of an incident support]
\lean{FABL.SixDisjointBlockChoice.incidentSupport_card}
\label{FABL.SixDisjointBlockChoice.incidentSupport_card}
\leanok
\uses{FABL.SixDisjointBlockChoice, FABL.SixDisjointBlockChoice.block, FABL.SixDisjointBlockChoice.block_card, FABL.SixDisjointBlockChoice.block_pairwiseDisjoint, FABL.SixDisjointBlockChoice.incidentSupport, FABL.card_k4IncidentEdges, FABL.k4IncidentEdges}
\difficulty{4}
Let $n,d$ be natural numbers, choose six pairwise disjoint $d$-element blocks indexed by
the edges of the complete graph on four vertices, and fix one of its four vertices. The
union of the blocks indexed by the three edges incident to that vertex has cardinality
$3d$.
\end{theorem}

\begin{theorem}[Uniqueness of a block containing an element]
\lean{FABL.SixDisjointBlockChoice.eq_of_mem_blocks}
\label{FABL.SixDisjointBlockChoice.eq_of_mem_blocks}
\leanok
\uses{FABL.SixDisjointBlockChoice, FABL.SixDisjointBlockChoice.block, FABL.SixDisjointBlockChoice.block_pairwiseDisjoint}
\difficulty{3}
Let $n,d$ be natural numbers and choose six pairwise disjoint $d$-element blocks from
the $n$-element set. If an element belongs to the blocks indexed by $i$ and $j$, then
$i=j$.
\end{theorem}

\begin{theorem}[Symmetric difference of four sets]
\lean{FABL.tupleSymmDiff_finFour}
\label{FABL.tupleSymmDiff_finFour}
\leanok
\uses{FABL.tupleSymmDiff}
\difficulty{1}
For any set type with decidable equality and any four finite subsets $S_0,S_1,S_2,S_3$,
their tuple symmetric difference equals
$S_0\mathbin{\triangle}(S_1\mathbin{\triangle}(S_2\mathbin{\triangle}(S_3\mathbin{\triangle}\varnothing)))$.
\end{theorem}

\begin{theorem}[Parity of the incident supports]
\lean{FABL.SixDisjointBlockChoice.tupleSymmDiff_incidentSupport}
\label{FABL.SixDisjointBlockChoice.tupleSymmDiff_incidentSupport}
\leanok
\uses{FABL.SixDisjointBlockChoice, FABL.SixDisjointBlockChoice.block, FABL.SixDisjointBlockChoice.eq_of_mem_blocks, FABL.SixDisjointBlockChoice.incidentSupport, FABL.k4EdgeEnds, FABL.k4IncidentEdges, FABL.tupleSymmDiff, FABL.tupleSymmDiff_finFour}
\difficulty{5}
For any natural numbers $n$ and $d$, take six pairwise disjoint $d$-element subsets of
an $n$-element set, indexed by the edges of the complete graph on four vertices. For
each vertex, let its incident support be the union of the three subsets indexed by edges
containing that vertex. The symmetric difference of the four incident supports is empty.
\end{theorem}

\begin{theorem}[Block as an intersection of incident supports]
\lean{FABL.SixDisjointBlockChoice.block_eq_incidentSupport_inter}
\label{FABL.SixDisjointBlockChoice.block_eq_incidentSupport_inter}
\leanok
\uses{FABL.SixDisjointBlockChoice, FABL.SixDisjointBlockChoice.block, FABL.SixDisjointBlockChoice.eq_of_mem_blocks, FABL.SixDisjointBlockChoice.incidentSupport, FABL.k4EdgeEnds, FABL.k4IncidentEdges}
\difficulty{5}
Let $n,d$ be natural numbers and choose six pairwise disjoint $d$-element blocks indexed
by the six edges of the complete graph on four vertices. For any edge, its block is the
intersection of the incident supports of its two endpoints, where each vertex's incident
support is the union of the blocks indexed by edges incident to that vertex.
\end{theorem}

\begin{definition}[Level-$k$ Boolean polynomial]
\lean{FABL.levelKPolynomial}
\statementsource{Analysis-of-Boolean-Functions-by-Ryan-ODonnell}{pdf-ffbe9c7b3306-p028-b002}
\label{FABL.levelKPolynomial}
\leanok
\uses{FABL.monomial}
For natural numbers $n$ and $k$, define a real-valued function on the Boolean cube
$\\{-1,1\\}^n$ by summing the monomial associated with each $k$-element subset of the
$n$ coordinates.
\end{definition}

\begin{theorem}[Fourier coefficients of a level polynomial]
\lean{FABL.fourierCoeff_levelKPolynomial}
\label{FABL.fourierCoeff_levelKPolynomial}
\leanok
\uses{FABL.expect_monomial_mul, FABL.fourierCoeff, FABL.levelKPolynomial, FABL.monomial}
\difficulty{4}
Let $n$ and $k$ be nonnegative integers, and let $T$ be a finite set of coordinates in
$\\{0,\\ldots,n-1\\}$. On the $n$-dimensional Boolean cube, form the polynomial by
summing the coordinate monomials over all $k$-element coordinate sets. Its Fourier
coefficient at $T$ is $1$ if $|T|=k$, and is $0$ otherwise.
\end{theorem}

\begin{theorem}[Second moment of a level polynomial]
\lean{FABL.expect_levelKPolynomial_sq}
\label{FABL.expect_levelKPolynomial_sq}
\leanok
\uses{FABL.fourierCoeff_levelKPolynomial, FABL.levelKPolynomial, FABL.parseval, FABL.uniformInner}
\difficulty{5}
For nonnegative integers $n$ and $k$, let $L_k$ be the sum, over all $k$-element subsets
of $\\{0,\\ldots,n-1\\}$, of the corresponding sign monomials on the $n$-dimensional
sign cube. Under the uniform distribution on that cube, $\\mathbb
E[L_k(x)^2]=\\binom{n}{k}$.
\end{theorem}

\begin{definition}[Balanced parity tuple]
\lean{FABL.LevelKParityTuple}
\statementsource{Analysis-of-Boolean-Functions-by-Ryan-ODonnell}{pdf-ffbe9c7b3306-p277-b004}
\label{FABL.LevelKParityTuple}
\leanok
\uses{FABL.tupleSymmDiff}
For natural numbers $n$, $k$, and $q$, a level-$k$ parity tuple is a family of $q$
subsets of an $n$-element set such that each subset has cardinality $k$ and the
symmetric difference of the entire family is empty.
\end{definition}

\begin{definition}[Enumeration of a tuple component]
\lean{FABL.LevelKParityTuple.enumeration}
\label{FABL.LevelKParityTuple.enumeration}
\leanok
\uses{FABL.LevelKParityTuple}
Let $q$ finite sets each have size $k$, with their symmetric difference empty, and let
one of these sets be specified. The enumeration of that component is a bijection from a
$k$-element index set to the elements of the specified set.
\end{definition}

\begin{definition}[Partwise permutations]
\lean{FABL.ParityTupleReordering}
\label{FABL.ParityTupleReordering}
\leanok
For natural numbers $q$ and $k$, a partwise reordering is a choice, for each of the $q$
parts, of a permutation of a $k$-element index set.
\end{definition}

\begin{definition}[Assigned vertex labels]
\lean{FABL.LevelKParityTuple.assignedVertexLabel}
\statementsource{Analysis-of-Boolean-Functions-by-Ryan-ODonnell}{pdf-ffbe9c7b3306-p277-b004}
\label{FABL.LevelKParityTuple.assignedVertexLabel}
\leanok
\uses{FABL.BookPartiteVertex, FABL.LevelKParityTuple, FABL.LevelKParityTuple.enumeration, FABL.ParityTupleReordering}
Let $n,k,q$ be natural numbers, and let a tuple consist of $q$ subsets of an $n$-element
set, each of size $k$, whose symmetric difference is empty. For each part, choose a
permutation of the $k$ indices. The assigned label of the vertex at position $i$ in that
part is the element in position given by the permutation, under the bijection from the
$k$ indices to that part’s subset.
\end{definition}

\begin{theorem}[Injectivity within each part]
\lean{FABL.LevelKParityTuple.assignedVertexLabel_injective_on_part}
\label{FABL.LevelKParityTuple.assignedVertexLabel_injective_on_part}
\leanok
\uses{FABL.LevelKParityTuple, FABL.LevelKParityTuple.assignedVertexLabel, FABL.LevelKParityTuple.enumeration, FABL.ParityTupleReordering}
\difficulty{3}
Let $n,k,q$ be natural numbers, and let $S_0,\ldots,S_{q-1}$ be subsets of
$\{0,\ldots,n-1\}$ whose symmetric difference is empty and each of which has cardinality
$k$. Choose a bijection from the $k$ positions to the elements of each subset and a
permutation of the positions for each part. Within any fixed part, the resulting vertex
labels are injective as a function of the position.
\end{theorem}

\begin{definition}[Equivalence of assigned-label fibers]
\lean{FABL.LevelKParityTuple.assignedLabelFiberEquivPartFiber}
\label{FABL.LevelKParityTuple.assignedLabelFiberEquivPartFiber}
\leanok
\uses{FABL.BookPartiteVertex, FABL.LevelKParityTuple, FABL.LevelKParityTuple.assignedVertexLabel, FABL.LevelKParityTuple.enumeration, FABL.ParityTupleReordering}
Let $q$ finite sets, each of size $k$, have empty symmetric difference, and choose a
permutation of the $k$ indices for each part. For any element $x$ of the ambient
$n$-element set, the labels assigned to $x$ form a set in bijection with the parts whose
component contains $x$.
\end{definition}

\begin{theorem}[Even size of each label fiber]
\lean{FABL.LevelKParityTuple.assignedVertexLabel_even_fiber}
\label{FABL.LevelKParityTuple.assignedVertexLabel_even_fiber}
\leanok
\uses{FABL.BookPartiteVertex, FABL.LevelKParityTuple, FABL.LevelKParityTuple.assignedLabelFiberEquivPartFiber, FABL.LevelKParityTuple.assignedVertexLabel, FABL.ParityTupleReordering, FABL.even_countP_of_tupleSymmDiff_eq_empty, FABL.natCard_subtype_eq_countP_ofFn}
\difficulty{4}
Let $n,k,q$ be natural numbers, and let $S_0,\ldots,S_{q-1}$ be subsets of
$\{0,\ldots,n-1\}$ whose symmetric difference is empty and each of which has cardinality
$k$. Choose a bijection from the $k$ positions to the elements of each $S_j$, and
independently permute the positions in each part before assigning the corresponding
elements as labels to the vertices. For every $x\in\{0,\ldots,n-1\}$, the number of
vertices labeled $x$ is even.
\end{theorem}

\begin{theorem}[Compatible perfect matching]
\lean{FABL.LevelKParityTuple.exists_compatible_matching}
\label{FABL.LevelKParityTuple.exists_compatible_matching}
\leanok
\uses{FABL.GraphPartitePerfectMatching, FABL.LevelKParityTuple, FABL.LevelKParityTuple.assignedVertexLabel, FABL.LevelKParityTuple.assignedVertexLabel_even_fiber, FABL.LevelKParityTuple.assignedVertexLabel_injective_on_part, FABL.MatchingCompatibleLabeling, FABL.ParityTupleReordering, FABL.exists_graphPartitePerfectMatching_compatible}
\difficulty{1}
Let $n,k,q$ be natural numbers, and let $(S_i)_{i=1}^q$ be finite subsets of $in n$
whose symmetric difference is empty and each of which has cardinality $k$. Choose any
permutation of the $k$ positions for each part. In the complete $q$-partite graph with
$k$ vertices in each part, there is a perfect matching such that the two endpoints of
every matching edge receive the same label in $in n$, where the label of a vertex is
the element of its part’s set selected by the chosen permutation and an enumeration of
that set.
\end{theorem}

\begin{definition}[Labels in a part]
\lean{FABL.matchingLabelSet}
\statementsource{Analysis-of-Boolean-Functions-by-Ryan-ODonnell}{pdf-ffbe9c7b3306-p277-b004}
\label{FABL.matchingLabelSet}
\leanok
\uses{FABL.BookPartiteVertex}
Let $n,q,k$ be natural numbers, let each vertex in a collection of $q$ parts of size $k$
have a label in an $n$-element set, and fix one part. The set of labels in that part is
the set of labels attained by its $k$ positions.
\end{definition}

\begin{theorem}[Labels in each part]
\lean{FABL.LevelKParityTuple.matchingLabelSet_assignedVertexLabel}
\label{FABL.LevelKParityTuple.matchingLabelSet_assignedVertexLabel}
\leanok
\uses{FABL.LevelKParityTuple, FABL.LevelKParityTuple.assignedVertexLabel, FABL.LevelKParityTuple.enumeration, FABL.ParityTupleReordering, FABL.matchingLabelSet}
\difficulty{4}
Let $n,k,q$ be natural numbers, and let $(S_i)_{i=1}^q$ be finite subsets of $in n$,
each of cardinality $k$. Choose any permutation of the $k$ positions for each part, and
label each position by the corresponding element under an enumeration of that part’s
set. For every part, the set of labels assigned to its $k$ positions is exactly that
part’s set.
\end{theorem}

\begin{definition}[Perfect matching with compatible labels]
\lean{FABL.GraphMatchingLabeling}
\statementsource{Analysis-of-Boolean-Functions-by-Ryan-ODonnell}{pdf-ffbe9c7b3306-p277-b004}
\label{FABL.GraphMatchingLabeling}
\leanok
\uses{FABL.BookPartiteVertex, FABL.GraphPartitePerfectMatching, FABL.MatchingCompatibleLabeling}
Let $q,k,n$ be natural numbers. A graph matching labeling consists of a perfect matching
in the complete $q$-partite graph with $k$ vertices in each part, together with a
labeling of its vertices by an $n$-element set, such that the endpoints of every
matching edge have the same label.
\end{definition}

\begin{definition}[Selected compatible perfect matching]
\lean{FABL.LevelKParityTuple.chosenCompatibleMatching}
\label{FABL.LevelKParityTuple.chosenCompatibleMatching}
\leanok
\uses{FABL.GraphPartitePerfectMatching, FABL.LevelKParityTuple, FABL.LevelKParityTuple.exists_compatible_matching, FABL.ParityTupleReordering}
For a tuple of $q$ subsets of an $n$-element set, each of size $k$ and with empty
symmetric difference, together with a permutation of the $k$ indices for each part, the
selected compatible perfect matching is a perfect matching in the complete $q$-partite
graph with $k$ vertices in each part.
\end{definition}

\begin{theorem}[Compatibility of the chosen matching]
\lean{FABL.LevelKParityTuple.chosenCompatibleMatching_spec}
\label{FABL.LevelKParityTuple.chosenCompatibleMatching_spec}
\leanok
\uses{FABL.LevelKParityTuple, FABL.LevelKParityTuple.assignedVertexLabel, FABL.LevelKParityTuple.chosenCompatibleMatching, FABL.LevelKParityTuple.exists_compatible_matching, FABL.MatchingCompatibleLabeling, FABL.ParityTupleReordering}
\difficulty{1}
Let $n,k,q$ be natural numbers, and let $S_0,\ldots,S_{q-1}$ be subsets of
$\{0,\ldots,n-1\}$ whose symmetric difference is empty and each of which has cardinality
$k$. Choose a bijection from the $k$ positions to the elements of each subset and a
permutation of the positions in each part. In the complete multipartite graph with $q$
parts of size $k$, the chosen perfect matching is compatible with the resulting vertex
labeling: the two endpoints of every matching edge have the same label.
\end{theorem}

\begin{definition}[Parity tuple matching lift]
\lean{FABL.parityTupleMatchingLift}
\statementsource{Analysis-of-Boolean-Functions-by-Ryan-ODonnell}{pdf-ffbe9c7b3306-p277-b004}
\label{FABL.parityTupleMatchingLift}
\leanok
\uses{FABL.GraphMatchingLabeling, FABL.LevelKParityTuple, FABL.LevelKParityTuple.assignedVertexLabel, FABL.LevelKParityTuple.chosenCompatibleMatching, FABL.LevelKParityTuple.chosenCompatibleMatching_spec, FABL.ParityTupleReordering}
Let $n,k,q$ be natural numbers, let $q$ subsets of an $n$-element set each have size $k$
and have empty symmetric difference, and choose a permutation of the $k$ indices for
each part. The parity tuple matching lift assigns each vertex the element in its part’s
subset indexed by its permuted position, and pairs this labeling with a compatible
perfect matching in the complete $q$-partite graph with $k$ vertices in each part.
\end{definition}

\begin{theorem}[Injectivity of the matching lift]
\lean{FABL.parityTupleMatchingLift_injective}
\label{FABL.parityTupleMatchingLift_injective}
\leanok
\uses{FABL.GraphMatchingLabeling, FABL.LevelKParityTuple, FABL.LevelKParityTuple.assignedVertexLabel, FABL.LevelKParityTuple.enumeration, FABL.LevelKParityTuple.matchingLabelSet_assignedVertexLabel, FABL.ParityTupleReordering, FABL.matchingLabelSet, FABL.parityTupleMatchingLift}
\difficulty{5}
For natural numbers $n,k,q$, consider tuples of $q$ subsets of an $n$-element set, each
of size $k$, whose symmetric difference is empty, together with a permutation of the $k$
positions for each of the $q$ parts. The map that sends each such pair to its chosen
perfect matching in the complete multipartite graph, equipped with the associated
compatible vertex labeling, is injective.
\end{theorem}

\begin{theorem}[Tuple reorderings]
\lean{FABL.card_parityTupleReordering}
\label{FABL.card_parityTupleReordering}
\leanok
\uses{FABL.ParityTupleReordering}
\difficulty{3}
For natural numbers $q$ and $k$, the set of $q$-tuples of permutations of a $k$-element
set has cardinality $(k!)^q$.
\end{theorem}

\begin{definition}[Parity tuple Fourier weight]
\lean{FABL.parityTupleFourierWeight}
\statementsource{Analysis-of-Boolean-Functions-by-Ryan-ODonnell}{pdf-ffbe9c7b3306-p277-b004}
\label{FABL.parityTupleFourierWeight}
\leanok
\uses{FABL.LevelKParityTuple, FABL.fourierCoeff}
Let $n,k,q$ be natural numbers, let $f$ be a real-valued function on the $n$-dimensional
Boolean cube, and let a tuple consist of $q$ subsets of an $n$-element set, each of size
$k$, with empty symmetric difference. Its Fourier weight for $f$ is the product, over
the parts, of the absolute values of the Fourier coefficients of $f$ indexed by the
corresponding subsets.
\end{definition}

\begin{definition}[Matching labeling Fourier weight]
\lean{FABL.matchingLabelingFourierWeight}
\statementsource{Analysis-of-Boolean-Functions-by-Ryan-ODonnell}{pdf-ffbe9c7b3306-p277-b004}
\label{FABL.matchingLabelingFourierWeight}
\leanok
\uses{FABL.GraphMatchingLabeling, FABL.fourierCoeff, FABL.matchingLabelSet}
Let $n,k,q$ be natural numbers, let $f$ be a real-valued function on the $n$-dimensional
Boolean cube, and let a graph matching labeling consist of a perfect matching in the
complete $q$-partite graph with $k$ vertices per part and a compatible labeling by an
$n$-element set. Its Fourier weight for $f$ is the product, over the parts, of the
absolute values of the Fourier coefficients indexed by the sets of labels attained in
those parts.
\end{definition}

\begin{theorem}[Fourier weight preservation under the lift]
\lean{FABL.parityTupleMatchingLift_weight}
\label{FABL.parityTupleMatchingLift_weight}
\leanok
\uses{FABL.LevelKParityTuple, FABL.LevelKParityTuple.assignedVertexLabel, FABL.LevelKParityTuple.matchingLabelSet_assignedVertexLabel, FABL.ParityTupleReordering, FABL.fourierCoeff, FABL.matchingLabelSet, FABL.matchingLabelingFourierWeight, FABL.parityTupleFourierWeight, FABL.parityTupleMatchingLift}
\difficulty{3}
For natural numbers $n,k,q$, a function $f:\{-1,1\}^n\to\mathbb{R}$, and a pair $z$
consisting of a tuple of $q$ size-$k$ subsets with empty symmetric difference and a
permutation of the $k$ positions in each part, the product of the absolute Fourier
coefficients indexed by the labels of the lifted matching equals the product of the
absolute Fourier coefficients indexed by the original tuple of subsets.
\end{theorem}

\begin{theorem}[Comparison of parity-tuple and matching sums]
\lean{FABL.parityTuple_sum_le_matchingLabeling_sum}
\label{FABL.parityTuple_sum_le_matchingLabeling_sum}
\leanok
\uses{FABL.GraphMatchingLabeling, FABL.LevelKParityTuple, FABL.ParityTupleReordering, FABL.card_parityTupleReordering, FABL.matchingLabelingFourierWeight, FABL.parityTupleFourierWeight, FABL.parityTupleMatchingLift, FABL.parityTupleMatchingLift_injective, FABL.parityTupleMatchingLift_weight, FABL.sum_le_inv_mul_sum_of_injective_lifts}
\difficulty{4}
For natural numbers $n,k,q$ and a function $f:\{-1,1\}^n\to\mathbb{R}$, the sum of the
parity-tuple Fourier weights is at most $(k!)^{-q}$ times the sum of the
matching-labeling Fourier weights. Here the first sum ranges over all $q$-tuples of
size-$k$ subsets of an $n$-element set with empty symmetric difference, and the second
ranges over all perfect matchings of the complete $q$-partite graph with $k$ vertices in
each part, equipped with compatible labels in the $n$-element set.
\end{theorem}

\begin{definition}[Fourier homogeneity]
\lean{FABL.IsFourierHomogeneous}
\statementsource{Analysis-of-Boolean-Functions-by-Ryan-ODonnell}{
  pdf-ffbe9c7b3306-p028-b002,
  pdf-ffbe9c7b3306-p277-b004
}
\label{FABL.IsFourierHomogeneous}
\leanok
\uses{FABL.fourierCoeff}
Let $n$ and $k$ be natural numbers, and let $f$ be a real-valued function on the
$n$-dimensional Boolean cube. The function is Fourier homogeneous of degree $k$ if every
nonzero Fourier coefficient of $f$ is indexed by a subset of size $k$.
\end{definition}

\begin{theorem}[Fourier weights of parity tuples]
\lean{FABL.sum_abs_fourierCoeff_tupleSymmDiff_eq_parityTuple_sum}
\label{FABL.sum_abs_fourierCoeff_tupleSymmDiff_eq_parityTuple_sum}
\leanok
\uses{FABL.IsFourierHomogeneous, FABL.LevelKParityTuple, FABL.fourierCoeff, FABL.parityTupleFourierWeight, FABL.tupleSymmDiff}
\difficulty{5}
Let $n,k,q$ be natural numbers and let $f$ be a real-valued function on the
$n$-dimensional sign cube whose nonzero Fourier coefficients are all indexed by sets of
cardinality $k$. Sum, over all ordered $q$-tuples of coordinate sets whose symmetric
difference is empty, the absolute value of the product of their Fourier coefficients.
This sum equals the sum of the products of the absolute Fourier coefficients over all
such tuples in which every coordinate set has cardinality $k$.
\end{theorem}

\begin{theorem}[Even moment bound by compatible labelings]
\lean{FABL.homogeneous_evenMoment_le_matchingLabeling_sum}
\proofsource{Analysis-of-Boolean-Functions-by-Ryan-ODonnell}{
  pdf-ffbe9c7b3306-p277-b004,
  pdf-ffbe9c7b3306-p278-b001
}
\label{FABL.homogeneous_evenMoment_le_matchingLabeling_sum}
\leanok
\uses{FABL.GraphMatchingLabeling, FABL.IsFourierHomogeneous, FABL.LevelKParityTuple, FABL.expect_pow_le_sum_abs_fourierCoeff_of_tupleSymmDiff_empty, FABL.fourierCoeff, FABL.matchingLabelingFourierWeight, FABL.parityTupleFourierWeight, FABL.parityTuple_sum_le_matchingLabeling_sum, FABL.sum_abs_fourierCoeff_tupleSymmDiff_eq_parityTuple_sum, FABL.tupleSymmDiff}
\difficulty{3}
Let $n,r,k$ be nonnegative integers, and let $f$ be a real-valued function on the sign
cube of dimension $n$ whose nonzero Fourier coefficients are all at level $k$. Its
$2r$-th moment under the uniform measure is at most $(k!)^{-2r}$ times the sum, over all
perfect matchings of the complete $2r$-partite graph with $k$ vertices in each part
together with vertex labelings by $\{1,\ldots,n\}$ that agree across every matching
edge, of the product over the $2r$ parts of the absolute Fourier coefficient indexed by
the set of labels appearing in that part.
\end{theorem}

\begin{theorem}[Finite product average by last coordinate]
\lean{FABL.expect_fin_snoc}
\label{FABL.expect_fin_snoc}
\leanok
\uses{FABL.finSnocEquiv}
\difficulty{3}
Let $\Omega$ be a finite nonempty set, let $m\in\mathbb N$, and let
$f:\Omega^{m+1}\to\mathbb R$. The uniform average of $f$ over all $(m+1)$-tuples equals
the uniform average over the first $m$ entries of the uniform average over the last
entry.
\end{theorem}

\begin{theorem}[Factoring two terms from a product]
\lean{FABL.prod_eq_mul_mul_prod_erase_two}
\label{FABL.prod_eq_mul_mul_prod_erase_two}
\leanok
\difficulty{3}
Let $P$ be a finite set, let $a,b\in P$ be distinct, and let $F:P\to\mathbb R$. Then
$\prod_{p\in P}F(p)=F(a)F(b)\prod_{p\in P\setminus\{a,b\}}F(p)$.
\end{theorem}

\begin{theorem}[Uniform product expectation bound]
\lean{FABL.uniform_expect_prod_le_prod_sqrt_expect_sq_twoRegular}
\label{FABL.uniform_expect_prod_le_prod_sqrt_expect_sq_twoRegular}
\leanok
\uses{FABL.expect_fin_snoc, FABL.prod_eq_mul_mul_prod_erase_two, FABL.uniformInner, FABL.uniformInner_le_uniformLpNorm_mul, FABL.uniformLpNorm, FABL.uniformLpNorm_two_eq_sqrt_expect_sq}
\difficulty{7}
Let $\Omega$ be a finite nonempty set, let $P$ be a finite set, and let $m$ be a natural
number. For each $p\in P$, choose a set $D_p\subseteq\{0,\ldots,m-1\}$, and suppose each
coordinate belongs to exactly two of these sets. Let $g_p$ be a nonnegative real-valued
function of assignments $a:\{0,\ldots,m-1\}\to\Omega$ that depends only on the
coordinates in $D_p$. When the coordinates of $a$ are chosen independently and uniformly
from $\Omega$, the expectation of $\prod_{p\in P}g_p(a)$ is at most $\prod_{p\in
P}\sqrt{\mathbb E[g_p(a)^2]}$.
\end{theorem}

\begin{theorem}[Finite-coordinate product expectation bound]
\lean{FABL.uniform_expect_prod_le_prod_sqrt_expect_sq_twoRegular_finite}
\label{FABL.uniform_expect_prod_le_prod_sqrt_expect_sq_twoRegular_finite}
\leanok
\uses{FABL.uniform_expect_prod_le_prod_sqrt_expect_sq_twoRegular}
\difficulty{5}
Let $\Omega$, $P$, and $E$ be finite sets, with $\Omega$ nonempty, and assign to each
$p\in P$ a subset $D_p\subseteq E$. Suppose each element of $E$ belongs to exactly two
of the sets $D_p$. For each $p$, let $g_p:\Omega^E\to\mathbb R$ be nonnegative and
depend only on the coordinates in $D_p$. With assignments chosen uniformly from
$\Omega^E$, the expectation of $\prod_{p\in P}g_p$ is at most $\prod_{p\in
P}\sqrt{\mathbb E[g_p^2]}$.
\end{theorem}

\begin{theorem}[Matching edge equality criterion]
\lean{FABL.matching_eq_or_adj_of_toEdge_eq}
\label{FABL.matching_eq_or_adj_of_toEdge_eq}
\leanok
\difficulty{4}
Let $G$ be a simple graph, let $M$ be a subgraph of $G$ that is a matching, and let
$v,w$ be vertices of $M$. If the edges assigned to $v$ and $w$ by the matching are
equal, then $v=w$ or $v$ and $w$ are adjacent in $M$.
\end{theorem}

\begin{definition}[Matching edges meeting a part]
\lean{FABL.graphMatchingPartEdges}
\statementsource{Analysis-of-Boolean-Functions-by-Ryan-ODonnell}{pdf-ffbe9c7b3306-p277-b004}
\label{FABL.graphMatchingPartEdges}
\leanok
\uses{FABL.BookPartiteVertex, FABL.GraphPartitePerfectMatching}
Let $q,k$ be natural numbers, and let a perfect matching be given in the complete
$q$-partite graph with $k$ vertices in each part. For a specified part, the associated
edge set consists of the matching edges having an endpoint in that part.
\end{definition}

\begin{theorem}[Edges in a part of a perfect matching]
\lean{FABL.mem_graphMatchingPartEdges_iff}
\label{FABL.mem_graphMatchingPartEdges_iff}
\leanok
\uses{FABL.BookPartiteVertex, FABL.GraphPartitePerfectMatching, FABL.graphMatchingPartEdges}
\difficulty{1}
Let $q$ and $k$ be natural numbers, and let a perfect matching be given in the complete
$q$-partite graph with $k$ vertices in each part. For any part $p$ and any edge of the
matching, that edge belongs to the subfamily of matching edges incident to a vertex in
part $p$ if and only if it contains a vertex of the form $(p,r)$ for some
$r\\in\\{0,\\ldots,k-1\\}$.
\end{theorem}

\begin{definition}[Matching edge at a part position]
\lean{FABL.graphMatchingPartEdge}
\statementsource{Analysis-of-Boolean-Functions-by-Ryan-ODonnell}{pdf-ffbe9c7b3306-p277-b004}
\label{FABL.graphMatchingPartEdge}
\leanok
\uses{FABL.GraphPartitePerfectMatching}
Let $q$ and $k$ be natural numbers, and let a perfect matching be given in the complete
multipartite graph with $q$ parts of size $k$. For a part $i in bbr q$ and a position
$j in bbr k$, the matching edge selected at the vertex $(i,j)$ is an edge of the
matching.
\end{definition}

\begin{theorem}[Matching edges incident to a part]
\lean{FABL.graphMatchingPartEdge_mem}
\label{FABL.graphMatchingPartEdge_mem}
\leanok
\uses{FABL.GraphPartitePerfectMatching, FABL.graphMatchingPartEdge, FABL.graphMatchingPartEdges, FABL.mem_graphMatchingPartEdges_iff}
\difficulty{3}
Let $q,k\in\mathbb{N}$, choose a perfect matching in the complete $q$-partite graph with
$k$ vertices per part, and fix a part and a position in it. The matching edge containing
the vertex at that position belongs to the set of matching edges incident to at least
one vertex in the fixed part.
\end{theorem}

\begin{theorem}[Injectivity of the partwise edge map]
\lean{FABL.graphMatchingPartEdge_injective}
\label{FABL.graphMatchingPartEdge_injective}
\leanok
\uses{FABL.GraphPartitePerfectMatching, FABL.graphMatchingPartEdge, FABL.matching_eq_or_adj_of_toEdge_eq}
\difficulty{4}
Let $q,k\in\mathbb{N}$, choose a perfect matching in the complete $q$-partite graph with
$k$ vertices per part, and fix a part. The map sending each position in that part to the
matching edge containing the corresponding vertex is injective.
\end{theorem}

\begin{definition}[Matching edges incident to a part]
\lean{FABL.graphMatchingPartEdgeEquiv}
\label{FABL.graphMatchingPartEdgeEquiv}
\leanok
\uses{FABL.GraphPartitePerfectMatching, FABL.graphMatchingPartEdge, FABL.graphMatchingPartEdge_injective, FABL.graphMatchingPartEdge_mem, FABL.graphMatchingPartEdges, FABL.mem_graphMatchingPartEdges_iff}
Let a perfect matching be given in the complete $q$-partite graph with $k$ vertices in
each part, and fix one part. The edges of the matching incident to a vertex in that part
are in bijection with a $k$-element index set.
\end{definition}

\begin{theorem}[Two parts incident to each matching edge]
\lean{FABL.graphMatchingPartEdges_degree_two}
\label{FABL.graphMatchingPartEdges_degree_two}
\leanok
\uses{FABL.GraphPartitePerfectMatching, FABL.graphMatchingPartEdges, FABL.mem_graphMatchingPartEdges_iff}
\difficulty{5}
Let $q,k\in\mathbb{N}$, choose a perfect matching in the complete $q$-partite graph with
$k$ vertices per part, and take any edge of the matching. Exactly two parts have a
matching edge incident to one of their vertices at some position; equivalently, the set
of such parts has cardinality $2$.
\end{theorem}

\begin{theorem}[Expectation reduced to relevant coordinates]
\lean{FABL.expect_eq_expect_subtype_of_dependsOn}
\label{FABL.expect_eq_expect_subtype_of_dependsOn}
\leanok
\difficulty{5}
Let $E$ and $\Omega$ be finite sets, with $\Omega$ nonempty. Let $D\subseteq E$ and let
$h:\Omega^E\to\mathbb{R}$ depend only on the coordinates in $D$. Then the uniform
average of $h$ over assignments on $E$ equals its uniform average over assignments on
$D$, extending each such assignment to $E$ by a fixed arbitrary element of $\Omega$
outside $D$.
\end{theorem}

\begin{definition}[Subsets of a fixed size]
\lean{FABL.LevelKSubset}
\label{FABL.LevelKSubset}
\leanok
For natural numbers $n$ and $k$, a level-$k$ subset of an $n$-element set is a subset
with exactly $k$ elements.
\end{definition}

\begin{definition}[Injective label tuple]
\lean{FABL.InjectiveLabelTuple}
\statementsource{Analysis-of-Boolean-Functions-by-Ryan-ODonnell}{pdf-ffbe9c7b3306-p277-b004}
\label{FABL.InjectiveLabelTuple}
\leanok
For natural numbers $n$ and $k$, an injective label tuple is an injective map from the
$k$ positions to the $n$ labels.
\end{definition}

\begin{definition}[Image of an injective tuple]
\lean{FABL.InjectiveLabelTuple.image}
\label{FABL.InjectiveLabelTuple.image}
\leanok
\uses{FABL.InjectiveLabelTuple, FABL.LevelKSubset}
Let $n$ and $k$ be natural numbers, and let $a:	ext{Fin }k	o	ext{Fin }n$ be injective.
The image of $a$ is the $k$-element subset $igl|	ext{range}(a)igr|$ of $	ext{Fin }n$.
\end{definition}

\begin{definition}[Indexing the image of an injective tuple]
\lean{FABL.InjectiveLabelTuple.imageEquiv}
\label{FABL.InjectiveLabelTuple.imageEquiv}
\leanok
\uses{FABL.InjectiveLabelTuple, FABL.InjectiveLabelTuple.image}
Let $n$ and $k$ be natural numbers, and let $a:	ext{Fin }k	o	ext{Fin }n$ be injective.
The map sending each index $i$ to $a(i)$ is a bijection from $	ext{Fin }k$ to the
elements of the image of $a$.
\end{definition}

\begin{definition}[Target for injective label codes]
\lean{FABL.InjectiveLabelCodeTarget}
\label{FABL.InjectiveLabelCodeTarget}
\leanok
\uses{FABL.LevelKSubset}
For natural numbers $n$ and $k$, let the target consist of pairs $(S,e)$, where $S$ is a
$k$-element subset of $	ext{Fin }n$ and $e$ is a bijection from $	ext{Fin }k$ to the
elements of $S$.
\end{definition}

\begin{definition}[Code of an injective tuple]
\lean{FABL.InjectiveLabelTuple.code}
\label{FABL.InjectiveLabelTuple.code}
\leanok
\uses{FABL.InjectiveLabelCodeTarget, FABL.InjectiveLabelTuple, FABL.InjectiveLabelTuple.image, FABL.InjectiveLabelTuple.imageEquiv}
Let $n$ and $k$ be natural numbers, and let $a:	ext{Fin }k	o	ext{Fin }n$ be injective.
Its code is the pair consisting of its image and the bijection from $	ext{Fin }k$ to
that image given by $i\mapsto a(i)$.
\end{definition}

\begin{theorem}[Injectivity of the tuple code]
\lean{FABL.injectiveLabelTuple_code_injective}
\label{FABL.injectiveLabelTuple_code_injective}
\leanok
\uses{FABL.InjectiveLabelCodeTarget, FABL.InjectiveLabelTuple, FABL.InjectiveLabelTuple.code}
\difficulty{3}
For nonnegative integers $n,k$, an injective $k$-tuple of elements of an $n$-element set
is encoded by its image, a $k$-element subset, together with a bijection from the $k$
positions to that subset. This encoding is injective: two tuples with the same code are
equal.
\end{theorem}

\begin{theorem}[Weighted sum over labeled subsets]
\lean{FABL.sum_injectiveLabelCodeTarget}
\label{FABL.sum_injectiveLabelCodeTarget}
\leanok
\uses{FABL.InjectiveLabelCodeTarget, FABL.LevelKSubset}
\difficulty{3}
Let $n$ and $k$ be natural numbers, and let $w$ assign a real number to each subset of
the coordinate set of size $k$ (viewed as a function on all subsets). Sum $w(S)$ over
all pairs consisting of such a subset $S$ and a bijection from a $k$-element set onto
$S$. This sum equals $k!$ times the sum of $w(S)$ over all subsets $S$ of size $k$.
\end{theorem}

\begin{theorem}[Weighted count of injective tuples]
\lean{FABL.sum_injectiveLabelTuple_le_factorial_mul}
\label{FABL.sum_injectiveLabelTuple_le_factorial_mul}
\leanok
\uses{FABL.InjectiveLabelCodeTarget, FABL.InjectiveLabelTuple, FABL.InjectiveLabelTuple.code, FABL.InjectiveLabelTuple.image, FABL.LevelKSubset, FABL.injectiveLabelTuple_code_injective, FABL.sum_injectiveLabelCodeTarget}
\difficulty{4}
Let $n,k$ be nonnegative integers, and let $w$ assign a real weight to each subset of
$\\{0,\\ldots,n-1\\}$, with $w(S)\\ge 0$ for every subset $S$. The sum of the weights of
the image sets of all injective maps from $\\{0,\\ldots,k-1\\}$ to $\\{0,\\ldots,n-1\\}$
is at most $k!$ times the sum of the weights of all $k$-element subsets of
$\\{0,\\ldots,n-1\\}$.
\end{theorem}

\begin{theorem}[Fourier coefficient sum over label tuples]
\lean{FABL.sum_labelTuple_fourierCoeff_sq_le}
\label{FABL.sum_labelTuple_fourierCoeff_sq_le}
\leanok
\uses{FABL.InjectiveLabelTuple, FABL.IsFourierHomogeneous, FABL.fourierCoeff, FABL.fourierWeightAtLevel, FABL.sum_injectiveLabelTuple_le_factorial_mul}
\difficulty{5}
Let $n,k$ be nonnegative integers, and let $f:\{-1,1\}^n\to\mathbb R$ have all its
nonzero Fourier coefficients supported on subsets of size $k$. Then the sum, over all
maps $a:[k]\to[n]$, of the squared Fourier coefficient indexed by the image of $a$ is at
most $k!$ times the total squared Fourier weight of $f$ on subsets of size $k$.
\end{theorem}

\begin{theorem}[Fourier coefficient bound for random tuples]
\lean{FABL.expect_labelTuple_fourierCoeff_sq_le}
\label{FABL.expect_labelTuple_fourierCoeff_sq_le}
\leanok
\uses{FABL.IsFourierHomogeneous, FABL.fourierCoeff, FABL.fourierWeightAtLevel, FABL.sum_labelTuple_fourierCoeff_sq_le}
\difficulty{4}
Let $n$ and $k$ be natural numbers, and let $f$ be a real-valued function on the
$n$-dimensional sign cube whose nonzero Fourier coefficients are indexed only by subsets
of size $k$. Choose a tuple $a$ uniformly from all maps from a $k$-element set to an
$n$-element set, allowing repeated values, and let $A$ be its image. Then the expected
squared Fourier coefficient indexed by $A$ is at most $k!/n^k$ times the total Fourier
weight of $f$ at level $k$.
\end{theorem}

\begin{definition}[Edge labels in one part]
\lean{FABL.graphMatchingEdgeLabelSet}
\statementsource{Analysis-of-Boolean-Functions-by-Ryan-ODonnell}{pdf-ffbe9c7b3306-p277-b004}
\label{FABL.graphMatchingEdgeLabelSet}
\leanok
\uses{FABL.GraphPartitePerfectMatching, FABL.graphMatchingPartEdge}
Let $q$, $k$, and $n$ be natural numbers, let a perfect matching be given in the
complete multipartite graph with $q$ parts of size $k$, and assign an $n$-valued label
to each edge of the matching. For a part $i$, the set of labels on the matching edges
selected at the vertices in that part consists of the labels obtained at every position
among the $k$ positions.
\end{definition}

\begin{definition}[Fourier factor for matching edge labels]
\lean{FABL.graphMatchingEdgeFourierFactor}
\statementsource{Analysis-of-Boolean-Functions-by-Ryan-ODonnell}{pdf-ffbe9c7b3306-p277-b004}
\label{FABL.graphMatchingEdgeFourierFactor}
\leanok
\uses{FABL.GraphPartitePerfectMatching, FABL.fourierCoeff, FABL.graphMatchingEdgeLabelSet}
Let $q$, $k$, and $n$ be natural numbers. Given a real-valued function on the
$n$-dimensional Boolean cube, a perfect matching in the complete multipartite graph with
$q$ parts of size $k$, a part $i$, and an assignment of an $n$-valued label to each
matching edge, the corresponding Fourier factor is the absolute value of the Fourier
coefficient indexed by the set of labels on the matching edges selected at the vertices
in part $i$.
\end{definition}

\begin{theorem}[Nonnegativity of a matching Fourier factor]
\lean{FABL.graphMatchingEdgeFourierFactor_nonneg}
\label{FABL.graphMatchingEdgeFourierFactor_nonneg}
\leanok
\uses{FABL.GraphPartitePerfectMatching, FABL.graphMatchingEdgeFourierFactor}
\difficulty{1}
Let $f$ be a real-valued function on the $n$-dimensional sign cube, let a perfect
matching be given in the complete multipartite graph with $q$ parts of $k$ vertices
each, fix a part, and assign labels in $\{0,\ldots,n-1\}$ to the matching edges. The
Fourier factor for that part is nonnegative.
\end{theorem}

\begin{theorem}[Dependence on edges incident to a part]
\lean{FABL.graphMatchingEdgeFourierFactor_dependsOn}
\label{FABL.graphMatchingEdgeFourierFactor_dependsOn}
\leanok
\uses{FABL.GraphPartitePerfectMatching, FABL.fourierCoeff, FABL.graphMatchingEdgeFourierFactor, FABL.graphMatchingPartEdge_mem, FABL.graphMatchingPartEdges}
\difficulty{3}
Let $f$ be a real-valued function on the $n$-dimensional sign cube, let a perfect
matching be given in the complete multipartite graph with $q$ parts of $k$ vertices
each, and fix one part. The Fourier factor for that part depends only on the labels
assigned to matching edges incident to it: any two edge assignments that agree on those
edges give the same factor.
\end{theorem}

\begin{theorem}[Dependence of the squared Fourier factor]
\lean{FABL.graphMatchingEdgeFourierFactor_sq_dependsOn}
\label{FABL.graphMatchingEdgeFourierFactor_sq_dependsOn}
\leanok
\uses{FABL.GraphPartitePerfectMatching, FABL.graphMatchingEdgeFourierFactor, FABL.graphMatchingEdgeFourierFactor_dependsOn, FABL.graphMatchingPartEdges}
\difficulty{2}
Let $f$ be a real-valued function on the $n$-dimensional sign cube, let a perfect
matching be given in the complete multipartite graph with $q$ parts of $k$ vertices
each, and fix one part. The square of its Fourier factor depends only on the labels
assigned to matching edges incident to that part: any two edge assignments agreeing on
those edges give the same squared factor.
\end{theorem}

\begin{theorem}[Mean square of a matching Fourier factor]
\lean{FABL.expect_graphMatchingEdgeFourierFactor_sq_eq}
\label{FABL.expect_graphMatchingEdgeFourierFactor_sq_eq}
\leanok
\uses{FABL.GraphPartitePerfectMatching, FABL.expect_eq_expect_subtype_of_dependsOn, FABL.fourierCoeff, FABL.graphMatchingEdgeFourierFactor, FABL.graphMatchingEdgeFourierFactor_sq_dependsOn, FABL.graphMatchingPartEdgeEquiv, FABL.graphMatchingPartEdge_mem, FABL.graphMatchingPartEdges}
\difficulty{4}
Let $q,k,n$ be natural numbers with $n>0$, let $f$ be a real-valued function on the
$n$-dimensional sign cube, and let a perfect matching of the complete $q$-partite graph
with $k$ vertices in each part be given, together with a part of that graph. Assign each
matching edge independently and uniformly a label from $\\{1,\\ldots,n\\}$. For each
assignment, take the absolute value of the Fourier coefficient of $f$ indexed by the set
of labels on the matching edges incident with the specified part. The expected square of
this quantity equals the expected square of the Fourier coefficient indexed by the set
of values of a uniformly random map from $\\{1,\\ldots,k\\}$ to $\\{1,\\ldots,n\\}$.
\end{theorem}

\begin{theorem}[Bound on a matching Fourier factor]
\lean{FABL.expect_graphMatchingEdgeFourierFactor_sq_le}
\label{FABL.expect_graphMatchingEdgeFourierFactor_sq_le}
\leanok
\uses{FABL.GraphPartitePerfectMatching, FABL.IsFourierHomogeneous, FABL.expect_graphMatchingEdgeFourierFactor_sq_eq, FABL.expect_labelTuple_fourierCoeff_sq_le, FABL.fourierWeightAtLevel, FABL.graphMatchingEdgeFourierFactor}
\difficulty{1}
Let $q,k,n$ be natural numbers with $n>0$. Let $f$ be a real-valued function on the
$n$-dimensional sign cube whose nonzero Fourier coefficients are supported only on
subsets of size $k$. Given a perfect matching of the complete $q$-partite graph with $k$
vertices in each part and a specified part, assign each matching edge independently and
uniformly a label from $\\{1,\\ldots,n\\}$. The expected square of the absolute Fourier
coefficient indexed by the set of labels on the matching edges incident with that part
is at most $\\frac{k!}{n^k}\\sum_{|S|=k}\\widehat f(S)^2$, where $\\widehat f(S)$
denotes the Fourier coefficient of $f$ indexed by $S$.
\end{theorem}

\begin{theorem}[Fourier factors of a partite matching]
\lean{FABL.graphMatching_finner}
\proofsource{Analysis-of-Boolean-Functions-by-Ryan-ODonnell}{pdf-ffbe9c7b3306-p278-b001}
\label{FABL.graphMatching_finner}
\leanok
\uses{FABL.GraphPartitePerfectMatching, FABL.graphMatchingEdgeFourierFactor, FABL.graphMatchingEdgeFourierFactor_dependsOn, FABL.graphMatchingEdgeFourierFactor_nonneg, FABL.graphMatchingPartEdges, FABL.graphMatchingPartEdges_degree_two, FABL.uniform_expect_prod_le_prod_sqrt_expect_sq_twoRegular_finite}
\difficulty{1}
Let $q,k,n$ be nonnegative integers with $n>0$, let $f$ be any real-valued function on
the Boolean cube $\\{-1,1\\}^n$, and let $M$ be a perfect matching in the complete
$q$-partite graph with $k$ vertices in each part. Assign independently and uniformly a
label from $[n]=\{0,\ldots,n-1\}$ to each edge of $M$. For each part, form the set of
labels on its incident matching edges, and let its factor be the absolute value of the
Fourier coefficient of $f$ indexed by that set, where the Fourier coefficient at $S$ is
the uniform average of $f(x)\prod_{i\in S}x_i$. The expected product of the part factors
is at most the product, over all parts, of the square root of the expected squared
factor for that part.
\end{theorem}

\begin{theorem}[Fixed matching Fourier bound]
\lean{FABL.graphMatching_fixed_expect_le}
\proofsource{Analysis-of-Boolean-Functions-by-Ryan-ODonnell}{
  pdf-ffbe9c7b3306-p278-b001,
  pdf-ffbe9c7b3306-p277-b004
}
\label{FABL.graphMatching_fixed_expect_le}
\leanok
\uses{FABL.GraphPartitePerfectMatching, FABL.IsFourierHomogeneous, FABL.expect_graphMatchingEdgeFourierFactor_sq_le, FABL.fourierWeight, FABL.fourierWeightAtLevel, FABL.graphMatchingEdgeFourierFactor, FABL.graphMatching_finner}
\difficulty{4}
Let $r,k,n$ be nonnegative integers with $n>0$, let $f$ be a real-valued function on
$\{-1,1\}^n$ whose nonzero Fourier coefficients are all indexed by sets of cardinality
$k$, and let $M$ be a perfect matching in the complete $2r$-partite graph with $k$
vertices per part. Assign independently and uniformly a label from
$[n]=\{0,\ldots,n-1\}$ to each edge of $M$. For each part, take the absolute value of
the Fourier coefficient of $f$ indexed by the set of labels on its incident matching
edges. The expected product of these factors is at most
$\left(\frac{k!}{n^k}W_k(f)\right)^r$, where $W_k(f)$ is the sum of the squares of the
Fourier coefficients indexed by $k$-element subsets of $[n]$.
\end{theorem}

\begin{theorem}[Matching edges incident to a part]
\lean{FABL.card_graphMatchingPartEdges}
\label{FABL.card_graphMatchingPartEdges}
\leanok
\uses{FABL.GraphPartitePerfectMatching, FABL.graphMatchingPartEdgeEquiv, FABL.graphMatchingPartEdges}
\difficulty{3}
Let $q,k$ be natural numbers, and let a perfect matching be given in the complete
$q$-partite graph with $k$ vertices in each part. For any one part, exactly $k$ matching
edges have an endpoint in that part.
\end{theorem}

\begin{theorem}[Size of a multipartite perfect matching]
\lean{FABL.card_graphMatching_edgeSet}
\label{FABL.card_graphMatching_edgeSet}
\leanok
\uses{FABL.GraphPartitePerfectMatching, FABL.card_graphMatchingPartEdges, FABL.graphMatchingPartEdges, FABL.graphMatchingPartEdges_degree_two}
\difficulty{5}
Let $r,k$ be natural numbers, and let a perfect matching be given in the complete
$2r$-partite graph with $k$ vertices in each part. The matching has exactly $rk$ edges.
\end{theorem}

\begin{theorem}[Fixed matching sum bound]
\lean{FABL.graphMatching_fixed_sum_le}
\proofsource{Analysis-of-Boolean-Functions-by-Ryan-ODonnell}{pdf-ffbe9c7b3306-p278-b001}
\label{FABL.graphMatching_fixed_sum_le}
\leanok
\uses{FABL.GraphPartitePerfectMatching, FABL.IsFourierHomogeneous, FABL.card_graphMatching_edgeSet, FABL.fourierWeightAtLevel, FABL.graphMatchingEdgeFourierFactor, FABL.graphMatching_fixed_expect_le}
\difficulty{4}
Let $r,k,n$ be nonnegative integers with $n>0$, let $f$ be a real-valued function on
$\{-1,1\}^n$ whose nonzero Fourier coefficients are all indexed by sets of cardinality
$k$, and let $M$ be a perfect matching in the complete $2r$-partite graph with $k$
vertices per part. Sum over all label assignments from $[n]=\{0,\ldots,n-1\}$ to the
edges of $M$ the product, over the parts, of the absolute Fourier coefficient of $f$
indexed by the set of labels on the part's incident matching edges. This sum is at most
$(k!)^r W_k(f)^r$, where $W_k(f)$ is the sum of the squares of the Fourier coefficients
indexed by $k$-element subsets of $[n]$.
\end{theorem}

\begin{definition}[Representative vertex of a matching edge]
\lean{FABL.graphMatchingEdgeRepresentative}
\label{FABL.graphMatchingEdgeRepresentative}
\leanok
\uses{FABL.GraphPartitePerfectMatching}
Let $q$ and $k$ be natural numbers, let a perfect matching be given in the complete
$q$-partite graph with $k$ vertices in each part, and let $e$ be one of its edges. The
chosen representative of $e$ is a vertex of the matching incident to $e$.
\end{definition}

\begin{theorem}[Representative of a matching edge]
\lean{FABL.graphMatchingEdgeRepresentative_spec}
\label{FABL.graphMatchingEdgeRepresentative_spec}
\leanok
\uses{FABL.GraphPartitePerfectMatching, FABL.graphMatchingEdgeRepresentative}
\difficulty{1}
Let $q,k\in\mathbb{N}$, and consider the complete $q$-partite graph whose vertices are
pairs $(i,j)$ with $0\le i<q$ and $0\le j<k$. For a perfect matching in this graph and
any edge of that matching, the chosen representative vertex maps to that same edge under
the matching's vertex-to-edge correspondence.
\end{theorem}

\begin{definition}[Vertex labels from edge assignments]
\lean{FABL.graphMatchingVertexLabel}
\label{FABL.graphMatchingVertexLabel}
\leanok
\uses{FABL.BookPartiteVertex, FABL.GraphPartitePerfectMatching}
Let $q,k,n$ be natural numbers, let a perfect matching be given in the complete
$q$-partite graph with $k$ vertices in each part, and assign a label in $	ext{Fin }n$ to
each matching edge. The induced labeling of the vertices assigns to each vertex the
label of the matching edge containing it.
\end{definition}

\begin{theorem}[Compatibility of vertex labels]
\lean{FABL.graphMatchingVertexLabel_compatible}
\label{FABL.graphMatchingVertexLabel_compatible}
\leanok
\uses{FABL.GraphPartitePerfectMatching, FABL.MatchingCompatibleLabeling, FABL.graphMatchingVertexLabel}
\difficulty{3}
Let $q,k,n\in\mathbb{N}$, choose a perfect matching in the complete $q$-partite graph
with $k$ vertices per part, and assign a label in $\{0,\ldots,n-1\}$ to each matching
edge. Give each vertex the label of its matching edge. Whenever two vertices are joined
by an edge of the chosen matching, their labels are equal.
\end{theorem}

\begin{definition}[Perfect matching with edge assignment]
\lean{FABL.GraphMatchingEdgeAssignment}
\statementsource{Analysis-of-Boolean-Functions-by-Ryan-ODonnell}{pdf-ffbe9c7b3306-p277-b004}
\label{FABL.GraphMatchingEdgeAssignment}
\leanok
\uses{FABL.GraphPartitePerfectMatching}
For natural numbers $q$, $k$, and $n$, a perfect matching with edge assignment consists
of a perfect matching in the complete multipartite graph with $q$ parts of size $k$,
together with an assignment of an $n$-valued label to each edge of that matching.
\end{definition}

\begin{definition}[Edge assignments induce compatible labelings]
\lean{FABL.graphMatchingEdgeAssignmentToLabeling}
\label{FABL.graphMatchingEdgeAssignmentToLabeling}
\leanok
\uses{FABL.GraphMatchingEdgeAssignment, FABL.GraphMatchingLabeling, FABL.graphMatchingVertexLabel, FABL.graphMatchingVertexLabel_compatible}
Let $q,k,n$ be natural numbers. From a perfect matching in the complete $q$-partite
graph with $k$ vertices in each part and an assignment of labels in $	ext{Fin }n$ to its
edges, form the vertex labeling that gives each vertex the label of its matching edge.
The resulting pair of matching and vertex labeling has equal labels at the endpoints of
every matching edge.
\end{definition}

\begin{definition}[Compatible labelings induce edge assignments]
\lean{FABL.graphMatchingLabelingToEdgeAssignment}
\label{FABL.graphMatchingLabelingToEdgeAssignment}
\leanok
\uses{FABL.GraphMatchingEdgeAssignment, FABL.GraphMatchingLabeling, FABL.graphMatchingEdgeRepresentative}
Let $q,k,n$ be natural numbers, and let a perfect matching in the complete $q$-partite
graph with $k$ vertices in each part be equipped with a vertex labeling by $	ext{Fin }n$
that gives equal labels to the endpoints of every matching edge. The induced edge
assignment sends each matching edge to the label of its chosen representative vertex.
\end{definition}

\begin{theorem}[Recovery of edge assignments]
\lean{FABL.graphMatchingEdgeAssignmentToLabeling_leftInverse}
\label{FABL.graphMatchingEdgeAssignmentToLabeling_leftInverse}
\leanok
\uses{FABL.GraphMatchingEdgeAssignment, FABL.GraphMatchingLabeling, FABL.graphMatchingEdgeAssignmentToLabeling, FABL.graphMatchingEdgeRepresentative, FABL.graphMatchingEdgeRepresentative_spec, FABL.graphMatchingLabelingToEdgeAssignment, FABL.matching_eq_or_adj_of_toEdge_eq}
\difficulty{5}
For a complete multipartite graph with $q$ parts of $k$ vertices each, a perfect
matching, and a natural number $n$, start with an assignment of labels in
$\{0,\ldots,n-1\}$ to the matching edges. Convert it to a vertex labeling by assigning
each vertex the label of its matched edge, then recover an edge assignment from that
labeling by choosing an endpoint representative for each edge. The recovered edge
assignment is the original assignment.
\end{theorem}

\begin{theorem}[Recovering edge assignments from labelings]
\lean{FABL.graphMatchingLabelingToEdgeAssignment_leftInverse}
\label{FABL.graphMatchingLabelingToEdgeAssignment_leftInverse}
\leanok
\uses{FABL.GraphMatchingEdgeAssignment, FABL.GraphMatchingLabeling, FABL.graphMatchingEdgeAssignmentToLabeling, FABL.graphMatchingEdgeRepresentative, FABL.graphMatchingEdgeRepresentative_spec, FABL.graphMatchingLabelingToEdgeAssignment}
\difficulty{3}
Let $q,k,n\in\mathbb{N}$. For any perfect matching in the complete $q$-partite graph
with $k$ vertices per part, assign a label in $\{0,\ldots,n-1\}$ to each matching edge
and give both endpoints that edge's label. Converting this labeling back to an edge
assignment recovers the original matching and every edge label.
\end{theorem}

\begin{definition}[Matching labelings and edge assignments]
\lean{FABL.graphMatchingLabelingEquivEdgeAssignment}
\label{FABL.graphMatchingLabelingEquivEdgeAssignment}
\leanok
\uses{FABL.GraphMatchingEdgeAssignment, FABL.GraphMatchingLabeling, FABL.graphMatchingEdgeAssignmentToLabeling, FABL.graphMatchingEdgeAssignmentToLabeling_leftInverse, FABL.graphMatchingLabelingToEdgeAssignment, FABL.graphMatchingLabelingToEdgeAssignment_leftInverse}
For natural numbers $q,k,n$, there is an equivalence between pairs consisting of a
perfect matching in the complete $q$-partite graph with $k$ vertices in each part and a
labeling of all its vertices by $	ext{Fin}(n)$ that is constant across every matched
edge, and a perfect matching together with an assignment of a label in $	ext{Fin}(n)$ to
each edge of that matching.
\end{definition}

\begin{theorem}[Factorization of matching Fourier weight]
\lean{FABL.graphMatchingEdgeAssignmentToLabeling_weight}
\label{FABL.graphMatchingEdgeAssignmentToLabeling_weight}
\leanok
\uses{FABL.GraphMatchingEdgeAssignment, FABL.graphMatchingEdgeAssignmentToLabeling, FABL.graphMatchingEdgeFourierFactor, FABL.matchingLabelingFourierWeight}
\difficulty{3}
Let $f$ be a real-valued function on the $n$-dimensional sign cube, and let an edge
assignment label the edges of a perfect matching in the complete multipartite graph with
$q$ parts of $k$ vertices each by elements of $\{0,\ldots,n-1\}$. The Fourier weight of
the induced vertex labeling is the product, over all parts, of the absolute values of
the Fourier coefficients indexed by the sets of labels assigned to the matching edges
incident to each part.
\end{theorem}

\begin{theorem}[Matching labeling sum as an edge assignment sum]
\lean{FABL.sum_matchingLabelingFourierWeight_eq_edgeAssignment}
\label{FABL.sum_matchingLabelingFourierWeight_eq_edgeAssignment}
\leanok
\uses{FABL.GraphMatchingEdgeAssignment, FABL.GraphMatchingLabeling, FABL.graphMatchingEdgeAssignmentToLabeling_weight, FABL.graphMatchingEdgeFourierFactor, FABL.graphMatchingLabelingEquivEdgeAssignment, FABL.matchingLabelingFourierWeight}
\difficulty{4}
Let $q,k,n$ be nonnegative integers and let $f:\{-1,1\}^n\to\mathbb R$. Form the
complete $q$-partite graph with $k$ vertices in each part. Sum, over every perfect
matching in this graph and every labeling of its vertices by $[n]$ that assigns equal
labels to the endpoints of each matched edge, the product over the parts of the absolute
value of the Fourier coefficient of $f$ indexed by the set of labels in that part. This
sum equals the sum over every perfect matching and every assignment of labels in $[n]$
to its edges of the product over the parts of the absolute value of the Fourier
coefficient indexed by the set of labels assigned to edges incident with that part.
\end{theorem}

\begin{theorem}[Fourier weight bound for matching labelings]
\lean{FABL.sum_matchingLabelingFourierWeight_le}
\proofsource{Analysis-of-Boolean-Functions-by-Ryan-ODonnell}{
  pdf-ffbe9c7b3306-p277-b004,
  pdf-ffbe9c7b3306-p278-b001
}
\label{FABL.sum_matchingLabelingFourierWeight_le}
\leanok
\uses{FABL.GraphMatchingLabeling, FABL.GraphPartitePerfectMatching, FABL.IsFourierHomogeneous, FABL.fourierWeightAtLevel, FABL.graphMatchingEdgeFourierFactor, FABL.graphMatching_fixed_sum_le, FABL.matchingLabelingFourierWeight, FABL.sum_matchingLabelingFourierWeight_eq_edgeAssignment}
\difficulty{4}
Let $r,k,n$ be nonnegative integers with $n>0$, and let $f:\{-1,1\}^n\to\mathbb R$ have
all its nonzero Fourier coefficients supported on subsets of size $k$. In the complete
$2r$-partite graph with $k$ vertices in each part, sum over all perfect matchings and
vertex labelings by $[n]$ that are constant on each matched edge the product, over the
parts, of the absolute Fourier coefficient indexed by the set of labels in that part.
This sum is at most the number of perfect matchings times $(k!)^r$ times the $r$th power
of the total squared Fourier weight on subsets of size $k$.
\end{theorem}

\begin{theorem}[Even moment bound by matching count]
\lean{FABL.homogeneous_evenMoment_le_graphMatching_count}
\proofsource{Analysis-of-Boolean-Functions-by-Ryan-ODonnell}{
  pdf-ffbe9c7b3306-p277-b004,
  pdf-ffbe9c7b3306-p278-b001
}
\label{FABL.homogeneous_evenMoment_le_graphMatching_count}
\leanok
\uses{FABL.GraphMatchingLabeling, FABL.GraphPartitePerfectMatching, FABL.IsFourierHomogeneous, FABL.fourierWeightAtLevel, FABL.homogeneous_evenMoment_le_matchingLabeling_sum, FABL.matchingLabelingFourierWeight, FABL.sum_matchingLabelingFourierWeight_le}
\difficulty{4}
Let $r,k,n$ be nonnegative integers with $n>0$, and let $f$ be a real-valued function on
the sign cube of dimension $n$ whose nonzero Fourier coefficients are all at level $k$.
Then its $2r$-th moment under the uniform measure is at most the number of perfect
matchings in the complete $2r$-partite graph with $k$ vertices in each part, divided by
$(k!)^r$, and multiplied by the $r$-th power of the total Fourier weight at level $k$.
\end{theorem}

\begin{theorem}[Proof heq]
\lean{FABL.proof_heq}
\label{FABL.proof_heq}
\leanok
% dataset-pipeline: stub -- prose pending build_informal_statements.py
Statement pending regeneration.
\end{theorem}

\begin{definition}[Endpoints of a pair]
\lean{FABL.pairEndpoints}
\label{FABL.pairEndpoints}
\leanok
For a pair $(a,b)$ of elements of a type, its endpoint list is the two-term list
$[a,b]$.
\end{definition}

\begin{definition}[Flattening a list of pairs]
\lean{FABL.flattenPairList}
\label{FABL.flattenPairList}
\leanok
\uses{FABL.pairEndpoints}
For a list of pairs of elements of a type, its flattened list is obtained by replacing
each pair with the list of its two entries and concatenating the resulting lists in
order.
\end{definition}

\begin{theorem}[Decoded matching endpoints form a permutation]
\lean{FABL.flattenPairList_decodePerfectMatchingCode_perm}
\label{FABL.flattenPairList_decodePerfectMatchingCode_perm}
\leanok
\uses{FABL.PerfectMatchingCode, FABL.decodePerfectMatchingCode, FABL.flattenPairList, FABL.pairEndpoints}
\difficulty{4}
Let $\alpha$ be a type, $m$ a natural number, $c$ a perfect-matching code of size $m$,
and $v$ a list of elements of $\alpha$ with length $2m$. The list obtained by decoding
$c$ from $v$ and then listing the endpoints of each decoded pair is a permutation of
$v$.
\end{theorem}

\begin{theorem}[Distinct endpoints in a decoded matching]
\lean{FABL.flattenPairList_decodePerfectMatchingCode_nodup}
\label{FABL.flattenPairList_decodePerfectMatchingCode_nodup}
\leanok
\uses{FABL.PerfectMatchingCode, FABL.decodePerfectMatchingCode, FABL.flattenPairList, FABL.flattenPairList_decodePerfectMatchingCode_perm}
\difficulty{1}
Let $\alpha$ be a type, $m$ a natural number, $c$ a perfect-matching code of size $m$,
and $v$ a list of elements of $\alpha$ with length $2m$ and no repetitions. The list
obtained by decoding $c$ from $v$ and then listing the endpoints of each decoded pair
has no repetitions.
\end{theorem}

\begin{definition}[Graph generated by a list of pairs]
\lean{FABL.pairListGraph}
\label{FABL.pairListGraph}
\leanok
Given a type of vertices and a list of ordered pairs of vertices, form the simple graph
whose edges are exactly the unordered pairs represented by entries of the list;
reversing an entry does not change its edge.
\end{definition}

\begin{theorem}[Adjacency in a graph from a pair list]
\lean{FABL.pairListGraph_adj_iff}
\label{FABL.pairListGraph_adj_iff}
\leanok
\uses{FABL.pairListGraph}
\difficulty{4}
Let $\alpha$ be a type, let $E$ be a list of ordered pairs from $\alpha$, and let
$u,v\in\alpha$. In the simple graph generated by the unordered pairs in $E$, vertices
$u$ and $v$ are adjacent if and only if either $(u,v)$ or $(v,u)$ occurs in $E$, and
$u\ne v$.
\end{theorem}

\begin{theorem}[A pair list defines a perfect matching]
\lean{FABL.pairListGraph_isPerfectMatching}
\label{FABL.pairListGraph_isPerfectMatching}
\leanok
\uses{FABL.flattenPairList, FABL.pairEndpoints, FABL.pairListGraph, FABL.pairListGraph_adj_iff}
\difficulty{6}
Let $\alpha$ be a type, let $V$ be a list of elements of $\alpha$, and let $E$ be a list
of pairs in $\alpha$. Suppose the list formed by concatenating the two endpoints of
every pair in $E$ is a permutation of $V$, that $V$ has no repetitions, and that every
element of $\alpha$ belongs to $V$. If $G$ is a simple graph on $\alpha$ containing the
graph generated by $E$ as a subgraph, then that generated graph, regarded as a subgraph
of $G$, is a perfect matching.
\end{theorem}

\begin{theorem}[Decoding commutes with relabelling]
\lean{FABL.decodePerfectMatchingCode_map}
\label{FABL.decodePerfectMatchingCode_map}
\leanok
\uses{FABL.PerfectMatchingCode, FABL.decodePerfectMatchingCode}
\difficulty{5}
Let $\alpha$ and $\beta$ be types, let $f:\alpha\to\beta$, let $m$ be a natural number,
and let $c$ be a perfect-matching code of order $m$. For every list $L$ of elements of
$\alpha$, decoding $c$ after applying $f$ to each element of $L$ gives the list obtained
by applying $f$ to both endpoints of each decoded pair:
$\operatorname{decode}(c,f(L))=\operatorname{map}((x,y)\mapsto(f(x),f(y)),\operatorname{decode}(c,L))$.
\end{theorem}

\begin{definition}[Enumeration of book vertices]
\lean{FABL.bookVertexFinEquiv}
\label{FABL.bookVertexFinEquiv}
\leanok
\uses{FABL.BookPartiteVertex}
For natural numbers $q$ and $k$, there is a bijection from the vertices of the complete
$q$-partite graph with $k$ vertices in each part to the integers from $0$ through
$qk-1$.
\end{definition}

\begin{definition}[List of finite vertices]
\lean{FABL.finVertexList}
\label{FABL.finVertexList}
\leanok
For a natural number $N$, the finite set $	ext{Fin}(N)$ is listed in its canonical
order.
\end{definition}

\begin{theorem}[Length of the finite vertex list]
\lean{FABL.length_finVertexList}
\label{FABL.length_finVertexList}
\leanok
\uses{FABL.finVertexList}
\difficulty{1}
For every natural number $N$, the list of vertices indexed by $\mathrm{Fin}(N)$ has
length $N$.
\end{theorem}

\begin{theorem}[Distinct finite indices]
\lean{FABL.nodup_finVertexList}
\label{FABL.nodup_finVertexList}
\leanok
\uses{FABL.finVertexList}
\difficulty{3}
For every natural number $N$, the list of all elements of $\text{Fin}(N)$ contains no
repetitions.
\end{theorem}

\begin{theorem}[Membership in the finite vertex list]
\lean{FABL.mem_finVertexList}
\label{FABL.mem_finVertexList}
\leanok
\uses{FABL.finVertexList}
\difficulty{3}
For every natural number $N$, every index in $	ext{Fin}(N)$ belongs to the list of all
indices in $	ext{Fin}(N)$.
\end{theorem}

\begin{theorem}[Values of the finite vertex list]
\lean{FABL.map_val_finVertexList}
\label{FABL.map_val_finVertexList}
\leanok
\uses{FABL.finVertexList}
\difficulty{4}
For every natural number $N$, applying the value map to each entry of the list of all
elements of $\{0,\ldots,N-1\}$ yields the list $0,1,\ldots,N-1$.
\end{theorem}

\begin{definition}[List of book vertices]
\lean{FABL.bookVertexList}
\label{FABL.bookVertexList}
\leanok
\uses{FABL.BookPartiteVertex, FABL.bookVertexFinEquiv, FABL.finVertexList}
For natural numbers $q$ and $k$, list the vertices of the complete $q$-partite graph
with $k$ vertices in each part by taking the canonical list of $	ext{Fin}(qk)$ and
transporting it back along the bijection between that set and the graph's vertices.
\end{definition}

\begin{theorem}[Length of the book vertex list]
\lean{FABL.length_bookVertexList}
\label{FABL.length_bookVertexList}
\leanok
\uses{FABL.bookVertexList}
\difficulty{1}
For all natural numbers $q$ and $k$, the book vertex list has length $qk$.
\end{theorem}

\begin{theorem}[Distinct book vertices]
\lean{FABL.nodup_bookVertexList}
\label{FABL.nodup_bookVertexList}
\leanok
\uses{FABL.bookVertexFinEquiv, FABL.bookVertexList, FABL.nodup_finVertexList}
\difficulty{1}
For any natural numbers $q$ and $k$, the list of vertices in the complete $q$-partite
indexing set with $k$ vertices in each part contains no repetitions.
\end{theorem}

\begin{theorem}[Membership in the book vertex list]
\lean{FABL.mem_bookVertexList}
\label{FABL.mem_bookVertexList}
\leanok
\uses{FABL.BookPartiteVertex, FABL.bookVertexFinEquiv, FABL.bookVertexList, FABL.mem_finVertexList}
\difficulty{3}
For natural numbers $q,k$, every vertex $(a,b)$ with $a\in\{0,\ldots,q-1\}$ and
$b\in\{0,\ldots,k-1\}$ belongs to the list of book vertices.
\end{theorem}

\begin{definition}[Decoded book-partite edges]
\lean{FABL.decodedBookPartiteEdges}
\label{FABL.decodedBookPartiteEdges}
\leanok
\uses{FABL.BookPartiteVertex, FABL.PerfectMatchingCode, FABL.bookVertexList, FABL.decodePerfectMatchingCode}
For natural numbers $q$ and $k$, and a code in the recursively defined type of
perfect-matching codes with parameter $(qk)/2$, the decoded book-partite edges are the
list obtained by applying the perfect-matching decoder to that code and the list of
book-partite vertices.
\end{definition}

\begin{theorem}[Index of a book vertex]
\lean{FABL.bookVertexFinEquiv_val_div}
\label{FABL.bookVertexFinEquiv_val_div}
\leanok
\uses{FABL.BookPartiteVertex, FABL.bookVertexFinEquiv}
\difficulty{3}
Let $q$ and $k$ be natural numbers with $k>0$. For a vertex $v$ in the book with $q$
parts and $k$ vertices per part, identify the vertices with the integers from $0$ to
$qk-1$ using the specified enumeration. Dividing the index of $v$ in this enumeration by
$k$ gives the index of its part.
\end{theorem}

\begin{theorem}[Decoded book edges under vertex coordinates]
\lean{FABL.decodedBookPartiteEdges_map_fin}
\label{FABL.decodedBookPartiteEdges_map_fin}
\leanok
\uses{FABL.PerfectMatchingCode, FABL.bookVertexFinEquiv, FABL.bookVertexList, FABL.decodePerfectMatchingCode, FABL.decodePerfectMatchingCode_map, FABL.decodedBookPartiteEdges, FABL.finVertexList}
\difficulty{4}
Let $q,k$ be natural numbers and let $c$ be a perfect-matching code of order $\lfloor
qk/2\rfloor$. Mapping both endpoints of every decoded book edge to their corresponding
elements of $\mathrm{Fin}(qk)$ gives exactly the list obtained by decoding $c$ from the
standard list of elements of $\mathrm{Fin}(qk)$.
\end{theorem}

\begin{theorem}[Decoded finite vertex list and its values]
\lean{FABL.decode_finVertexList_map_val}
\label{FABL.decode_finVertexList_map_val}
\leanok
\uses{FABL.PerfectMatchingCode, FABL.decodePerfectMatchingCode, FABL.decodePerfectMatchingCode_map, FABL.finVertexList, FABL.map_val_finVertexList}
\difficulty{3}
Let $m,N$ be natural numbers and let $c$ be a perfect-matching code of order $m$.
Decoding $c$ from the list of all elements of $\mathrm{Fin}(N)$ and taking the
natural-number values of both endpoints in each pair gives the same list of pairs as
decoding $c$ from $[0,1,\ldots,N-1]$.
\end{theorem}

\begin{theorem}[Decoded edges join distinct parts]
\lean{FABL.decodedBookPartiteEdge_parts_ne}
\label{FABL.decodedBookPartiteEdge_parts_ne}
\leanok
\uses{FABL.BookPartiteVertex, FABL.PartitePerfectMatching, FABL.bookVertexFinEquiv, FABL.bookVertexFinEquiv_val_div, FABL.decodePerfectMatchingCode, FABL.decode_finVertexList_map_val, FABL.decodedBookPartiteEdges, FABL.decodedBookPartiteEdges_map_fin, FABL.finVertexList}
\difficulty{5}
Let $q,k$ be natural numbers with $0<k$, and let $c$ be a partite perfect-matching code
for the book with $q$ parts of size $k$. If $(u,v)$ occurs among the edges decoded from
$c$, then the part containing $u$ differs from the part containing $v$.
\end{theorem}

\begin{theorem}[Decoded edges lie in the complete multipartite graph]
\lean{FABL.pairListGraph_decodedBookPartiteEdges_le}
\label{FABL.pairListGraph_decodedBookPartiteEdges_le}
\leanok
\uses{FABL.PartitePerfectMatching, FABL.bookPartiteGraph, FABL.decodedBookPartiteEdge_parts_ne, FABL.decodedBookPartiteEdges, FABL.pairListGraph, FABL.pairListGraph_adj_iff}
\difficulty{4}
Let $q,k$ be natural numbers with $k>0$, and let $c$ be a perfect-matching code of size
$\lfloor qk/2\rfloor$ whose decoded pairs join vertices from different parts of the
$q$-partite vertex set, with $k$ vertices in each part. The simple graph generated by
the decoded pairs is a subgraph of the complete $q$-partite graph on that vertex set.
\end{theorem}

\begin{definition}[Partite matching graph from a code]
\lean{FABL.partiteMatchingGraphOfCode}
\label{FABL.partiteMatchingGraphOfCode}
\leanok
\uses{FABL.GraphPartitePerfectMatching, FABL.PartitePerfectMatching, FABL.bookPartiteGraph, FABL.bookVertexList, FABL.decodedBookPartiteEdges, FABL.flattenPairList_decodePerfectMatchingCode_perm, FABL.length_bookVertexList, FABL.mem_bookVertexList, FABL.nodup_bookVertexList, FABL.pairListGraph, FABL.pairListGraph_decodedBookPartiteEdges_le, FABL.pairListGraph_isPerfectMatching}
Let $q,k$ be natural numbers with $k>0$ and $qk$ even. Every code representing a partite
perfect matching on the complete $q$-partite graph with $k$ vertices in each part
determines a perfect matching subgraph of that graph.
\end{definition}

\begin{theorem}[Injectivity of matching-code decoding]
\lean{FABL.pairListGraph_decodePerfectMatchingCode_injective}
\label{FABL.pairListGraph_decodePerfectMatchingCode_injective}
\leanok
\uses{FABL.PerfectMatchingCode, FABL.decodePerfectMatchingCode, FABL.flattenPairList, FABL.flattenPairList_decodePerfectMatchingCode_perm, FABL.pairEndpoints, FABL.pairListGraph, FABL.pairListGraph_adj_iff}
\difficulty{7}
Let $\alpha$ be a type, let $m$ be a natural number, and let $V$ be a list of elements
of $\alpha$ of length $2m$ with no repetitions. Decode each perfect-matching code of
size $m$ using $V$, and form the simple graph generated by the resulting list of pairs.
The map sending a code to this graph is injective.
\end{theorem}

\begin{definition}[Perfect matching on an even vertex set]
\lean{FABL.CompletePerfectMatching}
\statementsource{Analysis-of-Boolean-Functions-by-Ryan-ODonnell}{pdf-ffbe9c7b3306-p277-b004}
\label{FABL.CompletePerfectMatching}
\leanok
For a natural number $m$, a complete perfect matching consists of a perfect matching in
the complete graph on $2m$ vertices.
\end{definition}

\begin{definition}[Partner in a perfect matching]
\lean{FABL.perfectMatchingPartner}
\label{FABL.perfectMatchingPartner}
\leanok
\uses{FABL.CompletePerfectMatching}
For a natural number $m$, a perfect matching of the complete graph on $2m$ vertices, and
a vertex $v$ of that graph, define the partner of $v$ to be the vertex paired with $v$
by the matching.
\end{definition}

\begin{theorem}[Adjacency of matched partners]
\lean{FABL.perfectMatchingPartner_adj}
\label{FABL.perfectMatchingPartner_adj}
\leanok
\uses{FABL.CompletePerfectMatching, FABL.perfectMatchingPartner}
\difficulty{1}
For every natural number $m$, every perfect matching of the complete graph on $2m$
vertices, and every vertex $v$, the matched partner of $v$ is adjacent to $v$ in the
matching.
\end{theorem}

\begin{theorem}[Uniqueness of the matching partner]
\lean{FABL.perfectMatchingPartner_unique}
\label{FABL.perfectMatchingPartner_unique}
\leanok
\uses{FABL.CompletePerfectMatching, FABL.perfectMatchingPartner}
\difficulty{1}
Let $m$ be a nonnegative integer, let $M$ be a perfect matching of the complete graph on
$2m$ vertices, and let $v,w$ be vertices joined by an edge of $M$. Then $w$ is the
partner of $v$ in $M$.
\end{theorem}

\begin{theorem}[Involutive matching partner map]
\lean{FABL.perfectMatchingPartner_involutive}
\label{FABL.perfectMatchingPartner_involutive}
\leanok
\uses{FABL.CompletePerfectMatching, FABL.perfectMatchingPartner, FABL.perfectMatchingPartner_adj, FABL.perfectMatchingPartner_unique}
\difficulty{1}
Let $m$ be a nonnegative integer, and let $M$ be a perfect matching of the complete
graph on $2m$ vertices. The map sending each vertex to its unique partner in $M$ is an
involution: applying it twice returns the original vertex.
\end{theorem}

\begin{theorem}[Matching partners are distinct]
\lean{FABL.perfectMatchingPartner_ne}
\label{FABL.perfectMatchingPartner_ne}
\leanok
\uses{FABL.CompletePerfectMatching, FABL.perfectMatchingPartner, FABL.perfectMatchingPartner_adj}
\difficulty{3}
Let $m$ be a nonnegative integer, let $M$ be a perfect matching of the complete graph on
$2m$ vertices, and let $v$ be one of its vertices. The partner of $v$ in $M$ is
different from $v$.
\end{theorem}

\begin{definition}[Partner permutation]
\lean{FABL.perfectMatchingPartnerPerm}
\label{FABL.perfectMatchingPartnerPerm}
\leanok
\uses{FABL.CompletePerfectMatching, FABL.perfectMatchingPartner_involutive}
Let $m$ be a natural number and let a perfect matching be given on the complete graph
with vertex set of size $2m$. The permutation that sends each vertex to its partner in
the matching is an involution.
\end{definition}

\begin{theorem}[Matching partner permutation evaluation]
\lean{FABL.perfectMatchingPartnerPerm_apply}
\label{FABL.perfectMatchingPartnerPerm_apply}
\leanok
\uses{FABL.CompletePerfectMatching, FABL.perfectMatchingPartner, FABL.perfectMatchingPartnerPerm}
\difficulty{1}
For every natural number $m$, every perfect matching of the complete graph on $2m$
vertices, and every vertex $v$, the permutation associated with the matching sends $v$
to its matched partner.
\end{theorem}

\begin{theorem}[Injectivity of matching permutations]
\lean{FABL.perfectMatchingPartnerPerm_injective}
\label{FABL.perfectMatchingPartnerPerm_injective}
\leanok
\uses{FABL.CompletePerfectMatching, FABL.perfectMatchingPartner, FABL.perfectMatchingPartnerPerm, FABL.perfectMatchingPartner_adj, FABL.perfectMatchingPartner_unique}
\difficulty{5}
For every natural number $m$, distinct perfect matchings of the complete graph on $2m$
vertices yield distinct associated permutations.
\end{theorem}

\begin{definition}[Pairing permutations]
\lean{FABL.PairingPerm}
\label{FABL.PairingPerm}
\leanok
For a natural number $m$, a pairing permutation of $2m$ elements is a permutation of
$0,\ldots,2m-1$ whose cycle type consists of exactly $m$ cycles, each of length $2$.
\end{definition}

\begin{theorem}[Cycle type of the matching permutation]
\lean{FABL.perfectMatchingPartnerPerm_cycleType}
\label{FABL.perfectMatchingPartnerPerm_cycleType}
\leanok
\uses{FABL.CompletePerfectMatching, FABL.perfectMatchingPartner, FABL.perfectMatchingPartnerPerm, FABL.perfectMatchingPartner_involutive, FABL.perfectMatchingPartner_ne}
\difficulty{5}
For every natural number $m$ and every perfect matching of the complete graph on $2m$
vertices, the associated permutation has cycle type consisting of exactly $m$ cycles of
length $2$.
\end{theorem}

\begin{definition}[Permutation associated to a perfect matching]
\lean{FABL.completePerfectMatchingToPairingPerm}
\label{FABL.completePerfectMatchingToPairingPerm}
\leanok
\uses{FABL.CompletePerfectMatching, FABL.PairingPerm, FABL.perfectMatchingPartnerPerm, FABL.perfectMatchingPartnerPerm_cycleType}
Let $m$ be a natural number and let a perfect matching be given on the complete graph
with vertex set of size $2m$. Its partner permutation, viewed as a permutation of the
vertices, has cycle type consisting of exactly $m$ cycles, each of length $2$.
\end{definition}

\begin{theorem}[Injectivity of the matching-to-permutation map]
\lean{FABL.completePerfectMatchingToPairingPerm_injective}
\label{FABL.completePerfectMatchingToPairingPerm_injective}
\leanok
\uses{FABL.CompletePerfectMatching, FABL.PairingPerm, FABL.completePerfectMatchingToPairingPerm, FABL.perfectMatchingPartnerPerm_injective}
\difficulty{3}
For every natural number $m$, the map sending a perfect matching of the complete graph
on $2m$ vertices to the permutation that exchanges the two vertices in each matched pair
is injective. Its image consists of permutations whose cycle decomposition has exactly
$m$ cycles of length $2$.
\end{theorem}

\begin{theorem}[Number of pairings as permutations]
\lean{FABL.card_pairingPerm}
\label{FABL.card_pairingPerm}
\leanok
\uses{FABL.PairingPerm}
\difficulty{5}
For every natural number $m$, the number of permutations of $2m$ elements whose cycle
decomposition consists of exactly $m$ cycles of length two is $(2m-1)!!$.
\end{theorem}

\begin{definition}[Perfect matching graph from a code]
\lean{FABL.perfectMatchingGraphOfCode}
\label{FABL.perfectMatchingGraphOfCode}
\leanok
\uses{FABL.CompletePerfectMatching, FABL.PerfectMatchingCode, FABL.decodePerfectMatchingCode, FABL.finVertexList, FABL.flattenPairList_decodePerfectMatchingCode_perm, FABL.mem_finVertexList, FABL.nodup_finVertexList, FABL.pairListGraph, FABL.pairListGraph_isPerfectMatching}
Let $m$ be a natural number. Every code for a perfect matching on $2m$ vertices
determines a perfect matching subgraph of the complete graph on those vertices.
\end{definition}

\begin{theorem}[Injectivity of the matching graph map]
\lean{FABL.perfectMatchingGraphOfCode_injective}
\label{FABL.perfectMatchingGraphOfCode_injective}
\leanok
\uses{FABL.CompletePerfectMatching, FABL.PerfectMatchingCode, FABL.finVertexList, FABL.nodup_finVertexList, FABL.pairListGraph_decodePerfectMatchingCode_injective, FABL.perfectMatchingGraphOfCode}
\difficulty{4}
For every natural number $m$, the map sending a code in
$\mathrm{PerfectMatchingCode}(m)$ to its associated perfect matching in the complete
graph on $2m$ vertices is injective.
\end{theorem}

\begin{theorem}[Count of complete perfect matchings]
\lean{FABL.card_completePerfectMatching}
\label{FABL.card_completePerfectMatching}
\leanok
\uses{FABL.CompletePerfectMatching, FABL.PairingPerm, FABL.PerfectMatchingCode, FABL.card_pairingPerm, FABL.completePerfectMatchingToPairingPerm, FABL.completePerfectMatchingToPairingPerm_injective, FABL.natCard_perfectMatchingCode, FABL.perfectMatchingGraphOfCode, FABL.perfectMatchingGraphOfCode_injective}
\difficulty{4}
For every natural number $m$, the complete graph on $2m$ vertices has $(2m-1)!!$ perfect
matchings, where the count is understood as the natural-number double factorial
appearing in this formula.
\end{theorem}

\begin{theorem}[Perfect matching code bijection]
\lean{FABL.perfectMatchingGraphOfCode_bijective}
\label{FABL.perfectMatchingGraphOfCode_bijective}
\leanok
\uses{FABL.CompletePerfectMatching, FABL.PerfectMatchingCode, FABL.card_completePerfectMatching, FABL.natCard_perfectMatchingCode, FABL.perfectMatchingGraphOfCode, FABL.perfectMatchingGraphOfCode_injective}
\difficulty{3}
For every natural number $m$, the map that decodes a perfect matching code into a
perfect matching of the complete graph on $2m$ vertices is a bijection.
\end{theorem}

\begin{definition}[Matching codes and perfect matchings]
\lean{FABL.perfectMatchingCodeEquivCompletePerfectMatching}
\label{FABL.perfectMatchingCodeEquivCompletePerfectMatching}
\leanok
\uses{FABL.CompletePerfectMatching, FABL.PerfectMatchingCode, FABL.perfectMatchingGraphOfCode, FABL.perfectMatchingGraphOfCode_bijective}
For every natural number $m$, there is an equivalence between the codes for perfect
matchings of size $m$ and the perfect matchings of the complete graph on $2m$ vertices.
\end{definition}

\begin{theorem}[Injectivity of the code-to-matching map]
\lean{FABL.partiteMatchingGraphOfCode_injective}
\label{FABL.partiteMatchingGraphOfCode_injective}
\leanok
\uses{FABL.GraphPartitePerfectMatching, FABL.PartitePerfectMatching, FABL.bookVertexList, FABL.length_bookVertexList, FABL.nodup_bookVertexList, FABL.pairListGraph_decodePerfectMatchingCode_injective, FABL.partiteMatchingGraphOfCode}
\difficulty{4}
Let $q,k$ be natural numbers such that $k>0$ and $qk$ is even. The map sending a code
for a perfect matching on $qk$ vertices whose decoded edges join different parts to the
corresponding perfect matching of the complete $q$-partite graph with $k$ vertices in
each part is injective.
\end{theorem}

\begin{definition}[Even-sized partite vertex equivalence]
\lean{FABL.bookVertexEvenFinEquiv}
\label{FABL.bookVertexEvenFinEquiv}
\leanok
\uses{FABL.BookPartiteVertex, FABL.bookVertexFinEquiv}
For natural numbers $r$ and $k$, the vertices of the complete $(2r)$-partite graph with
$k$ vertices in each part are in bijection with the integers from $0$ through $2rk-1$.
\end{definition}

\begin{theorem}[Index of an even-partite book vertex]
\lean{FABL.bookVertexEvenFinEquiv_val_div}
\label{FABL.bookVertexEvenFinEquiv_val_div}
\leanok
\uses{FABL.BookPartiteVertex, FABL.bookVertexEvenFinEquiv, FABL.bookVertexFinEquiv, FABL.bookVertexFinEquiv_val_div}
\difficulty{3}
Let $r$ and $k$ be natural numbers with $k>0$. For a vertex $v$ in the book with $2r$
parts and $k$ vertices per part, identify the vertices with the integers from $0$ to
$2rk-1$ using the specified enumeration. Dividing the index of $v$ in this enumeration
by $k$ gives the index of its part.
\end{theorem}

\begin{definition}[Relabelled partite matching graph]
\lean{FABL.relabelBookMatchingGraph}
\label{FABL.relabelBookMatchingGraph}
\leanok
\uses{FABL.GraphPartitePerfectMatching, FABL.bookVertexEvenFinEquiv}
Let $r,k\in\mathbb{N}$, and let a perfect matching be given in the complete multipartite
graph with $2r$ parts of size $k$. Transport its edges to the vertex set of size $2rk$
using the canonical bijection from the partite vertices; the resulting graph has two
vertices adjacent exactly when their corresponding partite vertices are joined by the
matching.
\end{definition}

\begin{definition}[Partite matching as a complete matching]
\lean{FABL.graphPartiteMatchingToCompletePerfectMatching}
\label{FABL.graphPartiteMatchingToCompletePerfectMatching}
\leanok
\uses{FABL.CompletePerfectMatching, FABL.GraphPartitePerfectMatching, FABL.bookVertexEvenFinEquiv, FABL.relabelBookMatchingGraph}
Let $r,k\in\mathbb{N}$, and let a perfect matching be given in the complete multipartite
graph with $2r$ parts of size $k$. Relabel its vertices by the canonical bijection with
a set of size $2rk$; the resulting matching is a perfect matching in the complete graph
on those $2rk$ vertices.
\end{definition}

\begin{theorem}[Injectivity of the matching relabeling]
\lean{FABL.graphPartiteMatchingToCompletePerfectMatching_injective}
\label{FABL.graphPartiteMatchingToCompletePerfectMatching_injective}
\leanok
\uses{FABL.CompletePerfectMatching, FABL.GraphPartitePerfectMatching, FABL.bookVertexEvenFinEquiv, FABL.graphPartiteMatchingToCompletePerfectMatching, FABL.relabelBookMatchingGraph}
\difficulty{4}
For any nonnegative integers $r$ and $k$, the map sending a perfect matching in the
complete multipartite graph with $2r$ parts of size $k$ to the corresponding perfect
matching on $2rk$ vertices is injective.
\end{theorem}

\begin{theorem}[Decoded edges are matching graph edges]
\lean{FABL.decodedRangeEdge_perfectMatchingGraph_adj}
\label{FABL.decodedRangeEdge_perfectMatchingGraph_adj}
\leanok
\uses{FABL.PerfectMatchingCode, FABL.decodePerfectMatchingCode, FABL.decode_finVertexList_map_val, FABL.finVertexList, FABL.flattenPairList, FABL.flattenPairList_decodePerfectMatchingCode_nodup, FABL.flattenPairList_decodePerfectMatchingCode_perm, FABL.nodup_finVertexList, FABL.pairEndpoints, FABL.pairListGraph, FABL.pairListGraph_adj_iff, FABL.perfectMatchingGraphOfCode}
\difficulty{5}
Let $m$ be a natural number and let $c$ be a perfect-matching code of order $m$. If
$(a,b)$ occurs in the list decoded from $c$ and the list $[0,1,\ldots,2m-1]$, then
$a<2m$ and $b<2m$, and the vertices represented by $a$ and $b$ are adjacent in the
perfect-matching graph determined by $c$.
\end{theorem}

\begin{definition}[Even partite matching code condition]
\lean{FABL.IsEvenPartitePerfectMatchingCode}
\label{FABL.IsEvenPartitePerfectMatchingCode}
\leanok
\uses{FABL.PerfectMatchingCode, FABL.decodePerfectMatchingCode}
For $r,k\in\mathbb{N}$ and a matching code of size parameter $rk$, decode the code on
the list $0,1,\ldots,2rk-1$. The code is even partite precisely when every decoded pair
$(a,b)$ satisfies $\lfloor a/k\rfloor\ne\lfloor b/k\rfloor$.
\end{definition}

\begin{definition}[Even partite perfect matching]
\lean{FABL.EvenPartitePerfectMatching}
\label{FABL.EvenPartitePerfectMatching}
\leanok
\uses{FABL.IsEvenPartitePerfectMatchingCode, FABL.PerfectMatchingCode}
For $r,k\in\mathbb{N}$, an even partite perfect matching is a matching code with size
parameter $rk$ that satisfies the even partite matching condition.
\end{definition}

\begin{definition}[Equivalence of partite matching codes]
\lean{FABL.partitePerfectMatchingEvenEquiv}
\label{FABL.partitePerfectMatchingEvenEquiv}
\leanok
\uses{FABL.EvenPartitePerfectMatching, FABL.IsEvenPartitePerfectMatchingCode, FABL.IsPartitePerfectMatchingCode, FABL.PartitePerfectMatching}
For $r,k\in\mathbb{N}$, the codes for partite perfect matchings on $2r$ parts of size
$k$ are equivalent to the even partite perfect matching codes with parameters $r,k$.
\end{definition}

\begin{definition}[Raw code of a partite matching]
\lean{FABL.graphPartiteMatchingCodeRaw}
\label{FABL.graphPartiteMatchingCodeRaw}
\leanok
\uses{FABL.GraphPartitePerfectMatching, FABL.PerfectMatchingCode, FABL.graphPartiteMatchingToCompletePerfectMatching, FABL.perfectMatchingCodeEquivCompletePerfectMatching}
Let $r,k\in\mathbb{N}$, and let a perfect matching be given in the complete multipartite
graph with $2r$ parts of size $k$. Its raw matching code is the code corresponding to
the perfect matching obtained by relabelling its vertices as a set of size $2rk$.
\end{definition}

\begin{theorem}[Partite matching code property]
\lean{FABL.graphPartiteMatchingCodeRaw_isEvenPartite}
\label{FABL.graphPartiteMatchingCodeRaw_isEvenPartite}
\leanok
\uses{FABL.GraphPartitePerfectMatching, FABL.IsEvenPartitePerfectMatchingCode, FABL.bookVertexEvenFinEquiv, FABL.bookVertexEvenFinEquiv_val_div, FABL.decodedRangeEdge_perfectMatchingGraph_adj, FABL.graphPartiteMatchingCodeRaw, FABL.graphPartiteMatchingToCompletePerfectMatching, FABL.perfectMatchingCodeEquivCompletePerfectMatching, FABL.perfectMatchingGraphOfCode}
\difficulty{5}
Let $r,k$ be nonnegative integers with $k>0$, and let $M$ be a perfect matching in the
complete $2r$-partite graph with $k$ vertices in each part. Encode $M$ by the matching
code obtained from the canonical coding of perfect matchings on $2rk$ vertices. When
this code is decoded on the ordered list $0,1,\ldots,2rk-1$, every paired pair of
vertices has endpoints whose quotients upon division by $k$ are different.
\end{theorem}

\begin{definition}[Even code of a partite matching]
\lean{FABL.graphPartiteMatchingToEvenCode}
\label{FABL.graphPartiteMatchingToEvenCode}
\leanok
\uses{FABL.EvenPartitePerfectMatching, FABL.GraphPartitePerfectMatching, FABL.graphPartiteMatchingCodeRaw, FABL.graphPartiteMatchingCodeRaw_isEvenPartite}
Let $r,k\in\mathbb{N}$ with $k>0$, and let a perfect matching be given in the complete
multipartite graph with $2r$ parts of size $k$. The associated even code is its raw
matching code, together with the assertion that it satisfies the even partite matching
condition.
\end{definition}

\begin{theorem}[Injectivity of the even-partite code map]
\lean{FABL.graphPartiteMatchingToEvenCode_injective}
\label{FABL.graphPartiteMatchingToEvenCode_injective}
\leanok
\uses{FABL.EvenPartitePerfectMatching, FABL.GraphPartitePerfectMatching, FABL.graphPartiteMatchingToCompletePerfectMatching_injective, FABL.graphPartiteMatchingToEvenCode, FABL.perfectMatchingCodeEquivCompletePerfectMatching}
\difficulty{3}
Let $r$ and $k$ be nonnegative integers with $k>0$. The map sending a perfect matching
in the complete multipartite graph with $2r$ parts of size $k$ to its even-partite
perfect-matching code is injective.
\end{theorem}

\begin{theorem}[Matching count for positive part size]
\lean{FABL.card_graphPartitePerfectMatching_eq_partitePerfectMatching_of_pos}
\label{FABL.card_graphPartitePerfectMatching_eq_partitePerfectMatching_of_pos}
\leanok
\uses{FABL.EvenPartitePerfectMatching, FABL.GraphPartitePerfectMatching, FABL.PartitePerfectMatching, FABL.graphPartiteMatchingToEvenCode, FABL.graphPartiteMatchingToEvenCode_injective, FABL.partiteMatchingGraphOfCode, FABL.partiteMatchingGraphOfCode_injective, FABL.partitePerfectMatchingEvenEquiv}
\difficulty{5}
For natural numbers $r$ and $k$ with $k>0$, the number of perfect matchings in the
complete $2r$-partite graph with $k$ vertices in each part equals the number of perfect
matching codes of length parameter $(2rk)/2$ whose decoded pairs lie in distinct parts.
\end{theorem}

\begin{definition}[Perfect matching on the empty partite graph]
\lean{FABL.emptyGraphPartitePerfectMatching}
\label{FABL.emptyGraphPartitePerfectMatching}
\leanok
\uses{FABL.GraphPartitePerfectMatching}
For every $r\in\mathbb{N}$, the complete multipartite graph with $2r$ parts of size zero
has a perfect matching, namely the matching consisting of no edges on its empty vertex
set.
\end{definition}

\begin{theorem}[Matching count with empty parts]
\lean{FABL.card_graphPartitePerfectMatching_zero}
\label{FABL.card_graphPartitePerfectMatching_zero}
\leanok
\uses{FABL.CompletePerfectMatching, FABL.GraphPartitePerfectMatching, FABL.card_completePerfectMatching, FABL.emptyGraphPartitePerfectMatching, FABL.graphPartiteMatchingToCompletePerfectMatching, FABL.graphPartiteMatchingToCompletePerfectMatching_injective}
\difficulty{4}
For every natural number $r$, the complete $2r$-partite graph with zero vertices in each
part has exactly one perfect matching.
\end{theorem}

\begin{theorem}[Partite matchings on empty parts]
\lean{FABL.card_partitePerfectMatching_zero}
\label{FABL.card_partitePerfectMatching_zero}
\leanok
\uses{FABL.IsPartitePerfectMatchingCode, FABL.PartitePerfectMatching, FABL.PerfectMatchingCode, FABL.decodePerfectMatchingCode}
\difficulty{5}
For every natural number $r$, the set of perfect matchings on zero vertices divided into
$2r$ parts of size zero, with every matched pair in different parts, has cardinality
one.
\end{theorem}

\begin{theorem}[Equality of matching counts]
\lean{FABL.card_graphPartitePerfectMatching_eq_partitePerfectMatching}
\label{FABL.card_graphPartitePerfectMatching_eq_partitePerfectMatching}
\leanok
\uses{FABL.GraphPartitePerfectMatching, FABL.PartitePerfectMatching, FABL.card_graphPartitePerfectMatching_eq_partitePerfectMatching_of_pos, FABL.card_graphPartitePerfectMatching_zero, FABL.card_partitePerfectMatching_zero}
\difficulty{4}
For all natural numbers $r$ and $k$, the number of perfect matchings in the complete
$2r$-partite graph with $k$ vertices in each part equals the number of perfect matching
codes of length parameter $(2rk)/2$ whose decoded pairs lie in distinct parts.
\end{theorem}

\begin{definition}[Partite matchings and codes]
\lean{FABL.graphPartitePerfectMatchingEquivCode}
\label{FABL.graphPartitePerfectMatchingEquivCode}
\leanok
\uses{FABL.GraphPartitePerfectMatching, FABL.PartitePerfectMatching, FABL.card_graphPartitePerfectMatching_eq_partitePerfectMatching}
For natural numbers $r$ and $k$, consider the complete multipartite graph with $2r$
parts, each containing $k$ vertices. There is an equivalence between its perfect
matchings and recursive pairing codes of length parameter $(2rk)/2$ whose decoded pairs
join vertices from different parts.
\end{definition}

\begin{theorem}[Even moment bound for homogeneous functions]
\lean{FABL.homogeneous_evenMoment_le_partiteMatching_count}
\proofsource{Analysis-of-Boolean-Functions-by-Ryan-ODonnell}{
  pdf-ffbe9c7b3306-p277-b004,
  pdf-ffbe9c7b3306-p278-b001
}
\label{FABL.homogeneous_evenMoment_le_partiteMatching_count}
\leanok
\uses{FABL.IsFourierHomogeneous, FABL.PartitePerfectMatching, FABL.card_graphPartitePerfectMatching_eq_partitePerfectMatching, FABL.fourierWeightAtLevel, FABL.homogeneous_evenMoment_le_graphMatching_count}
\difficulty{1}
Let $n>0$ and let $r,k$ be nonnegative integers. Let $f$ be a real-valued function on
the sign cube $\\{-1,1\\}^n$ whose nonzero Fourier coefficients are supported only on
sets of size $k$. Let $M_{r,k}$ be the finite set of codes of length $rk$ whose decoded
list of pairs on the labels $0,\ldots,2rk-1$ has every pair joining labels in different
consecutive blocks of $k$ labels. Then
\[\mathbb E_x f(x)^{2r}\leq \frac{|M_{r,k}|}{(k!)^r}\,W_k(f)^r,\]
where $W_k(f)$ is the sum of the squared Fourier coefficients indexed by sets of size
$k$.
\end{theorem}

\begin{theorem}[Fourier weight of a homogeneous function]
\lean{FABL.fourierWeightAtLevel_eq_uniformLpNorm_two_sq_of_homogeneous}
\label{FABL.fourierWeightAtLevel_eq_uniformLpNorm_two_sq_of_homogeneous}
\leanok
\uses{FABL.IsFourierHomogeneous, FABL.fourierCoeff, FABL.fourierWeight, FABL.fourierWeightAtLevel, FABL.parseval, FABL.uniformLpNorm, FABL.uniformLpNorm_two_sq_eq_uniformInner}
\difficulty{4}
Let $f:\{-1,1\}^n\to\mathbb R$ have all its nonzero Fourier coefficients supported on
subsets of cardinality $k$, where $k$ is a natural number. Then the total Fourier weight
of $f$ at level $k$ equals the square of its uniform $L^2$ norm.
\end{theorem}

\begin{theorem}[Homogeneous uniform lp norm le partite matching]
\lean{FABL.homogeneous_uniformLpNorm_le_partiteMatching}
\statementsource{Analysis-of-Boolean-Functions-by-Ryan-ODonnell}{
  pdf-ffbe9c7b3306-p277-b004,
  pdf-ffbe9c7b3306-p278-b001
}
\label{FABL.homogeneous_uniformLpNorm_le_partiteMatching}
\leanok
\uses{FABL.IsFourierHomogeneous, FABL.PartitePerfectMatching, FABL.fourierWeightAtLevel, FABL.fourierWeightAtLevel_eq_uniformLpNorm_two_sq_of_homogeneous, FABL.homogeneous_evenMoment_le_partiteMatching_count, FABL.uniformLpNorm, FABL.uniformLpNorm_even_pow_eq_expect_even, FABL.uniformLpNorm_nonneg}
\difficulty{6}
% dataset-pipeline: stub -- prose pending build_informal_statements.py
Statement pending regeneration.
\end{theorem}

\begin{definition}[Vertices in a part]
\lean{FABL.remainingPartSlice}
\label{FABL.remainingPartSlice}
\leanok
\uses{FABL.BookPartiteVertex}
Given natural numbers $r,k$, a finite set of vertices in the $2r$-partite vertex set,
and one of the $2r$ parts, the corresponding part slice is the subset of those vertices
lying in that part.
\end{definition}

\begin{definition}[Largest remaining part]
\lean{FABL.greedyLargestPart}
\label{FABL.greedyLargestPart}
\leanok
\uses{FABL.BookPartiteVertex, FABL.remainingPartSlice}
Let $r,k$ be natural numbers with $r>0$, and let a finite set of vertices be given in a
$2r$-partite vertex set. The selected part is one whose slice of the given set has
maximum cardinality among all parts.
\end{definition}

\begin{theorem}[Largest part slice]
\lean{FABL.remainingPartSlice_card_le_greedyLargestPart}
\label{FABL.remainingPartSlice_card_le_greedyLargestPart}
\leanok
\uses{FABL.BookPartiteVertex, FABL.greedyLargestPart, FABL.remainingPartSlice}
\difficulty{3}
Let $r$ be a positive integer, let $k$ be a nonnegative integer, and let $R$ be a finite
set of vertices partitioned into $2r$ parts, each containing $k$ possible vertices. The
number of vertices of $R$ in any one part is at most the number in the part selected to
have largest such count.
\end{theorem}

\begin{theorem}[Counting vertices by part]
\lean{FABL.sum_card_remainingPartSlice}
\label{FABL.sum_card_remainingPartSlice}
\leanok
\uses{FABL.BookPartiteVertex, FABL.remainingPartSlice}
\difficulty{4}
Let $r$ and $k$ be nonnegative integers, and let $R$ be a finite set of vertices
partitioned into $2r$ parts, each with $k$ possible vertices. The sum, over all parts,
of the number of vertices of $R$ in that part equals the total number of vertices in
$R$.
\end{theorem}

\begin{theorem}[Bound by the largest part]
\lean{FABL.remaining_card_le_parts_mul_largest}
\label{FABL.remaining_card_le_parts_mul_largest}
\leanok
\uses{FABL.BookPartiteVertex, FABL.greedyLargestPart, FABL.remainingPartSlice, FABL.remainingPartSlice_card_le_greedyLargestPart, FABL.sum_card_remainingPartSlice}
\difficulty{4}
Let $r$ be a positive integer, let $k$ be a nonnegative integer, and let $R$ be a finite
set of vertices partitioned into $2r$ parts, each containing $k$ possible vertices. The
size of $R$ is at most $2r$ times the number of its vertices in a part having the
largest such count.
\end{theorem}

\begin{theorem}[Nonempty largest-part slice]
\lean{FABL.greedyLargestPartSlice_nonempty}
\label{FABL.greedyLargestPartSlice_nonempty}
\leanok
\uses{FABL.BookPartiteVertex, FABL.greedyLargestPart, FABL.remainingPartSlice, FABL.remainingPartSlice_card_le_greedyLargestPart}
\difficulty{4}
Let $r>0$, let $k$ be a nonnegative integer, and let $R$ be a nonempty finite set of
vertices from a complete multipartite graph with $2r$ parts of size $k$. The slice of
$R$ in a part of maximum slice size is nonempty.
\end{theorem}

\begin{definition}[Greedy pivot vertex]
\lean{FABL.greedyPivot}
\label{FABL.greedyPivot}
\leanok
\uses{FABL.BookPartiteVertex, FABL.greedyLargestPartSlice_nonempty}
Let $r,k$ be natural numbers with $r>0$, and let a nonempty finite set of vertices be
given in a $2r$-partite vertex set. The selected pivot is a vertex in a part of maximum
cardinality within that set.
\end{definition}

\begin{theorem}[Greedy pivot belongs to the remaining set]
\lean{FABL.greedyPivot_mem}
\label{FABL.greedyPivot_mem}
\leanok
\uses{FABL.BookPartiteVertex, FABL.greedyLargestPartSlice_nonempty, FABL.greedyPivot}
\difficulty{1}
Let $r>0$, let $k$ be a nonnegative integer, and let $R$ be a nonempty finite set of
vertices from a complete multipartite graph with $2r$ parts of size $k$. The pivot
chosen from a part of maximum slice size belongs to $R$.
\end{theorem}

\begin{theorem}[Part of the greedy pivot]
\lean{FABL.greedyPivot_part}
\label{FABL.greedyPivot_part}
\leanok
\uses{FABL.BookPartiteVertex, FABL.greedyLargestPart, FABL.greedyLargestPartSlice_nonempty, FABL.greedyPivot}
\difficulty{2}
Let $r>0$, let $k$ be a nonnegative integer, and let $R$ be a nonempty finite set of
vertices from a complete multipartite graph with $2r$ parts of size $k$. The pivot
chosen from a part of maximum slice size lies in that part.
\end{theorem}

\begin{definition}[Greedy partner choices]
\lean{FABL.greedyPartnerChoices}
\label{FABL.greedyPartnerChoices}
\leanok
\uses{FABL.BookPartiteVertex, FABL.greedyPivot}
Let $r,k$ be natural numbers with $r>0$, and let a nonempty finite set of vertices be
given in a $2r$-partite vertex set. The partner choices are the vertices in that set
whose part differs from the part containing the selected pivot.
\end{definition}

\begin{theorem}[Greedy choice count decomposition]
\lean{FABL.card_greedyPartnerChoices_add_largest}
\label{FABL.card_greedyPartnerChoices_add_largest}
\leanok
\uses{FABL.BookPartiteVertex, FABL.greedyLargestPart, FABL.greedyPartnerChoices, FABL.greedyPivot_part, FABL.remainingPartSlice}
\difficulty{3}
Let $r,k$ be natural numbers with $r>0$, and let $R$ be a nonempty finite set of
vertices $(a,j)$ with $a\\in\\{0,\\ldots,2r-1\\}$ and $j\\in\\{0,\\ldots,k-1\\}$. Choose
a part of largest size among the $2r$ parts, and let $P$ be the vertex chosen from that
part by the greedy pivot rule. The number of vertices of $R$ whose part differs from the
part of $P$, plus the number of vertices of $R$ in the chosen largest part, equals
$|R|$.
\end{theorem}

\begin{theorem}[Greedy partner choice bound]
\lean{FABL.card_greedyPartnerChoices_le}
\label{FABL.card_greedyPartnerChoices_le}
\leanok
\uses{FABL.BookPartiteVertex, FABL.card_greedyPartnerChoices_add_largest, FABL.greedyLargestPart, FABL.greedyPartnerChoices, FABL.remainingPartSlice, FABL.remaining_card_le_parts_mul_largest}
\difficulty{5}
Let $r,k,s$ be natural numbers with $r>0$ and $s>0$, and let $R$ be a finite set of
vertices in $2r$ parts of size $k$ with $|R|=2s$. Choose a part containing the largest
number of vertices of $R$, and choose a pivot vertex from that part. The number of
vertices of $R$ in parts other than the pivot's part is at most $\frac{2r-1}{r}s$.
\end{theorem}

\begin{definition}[Partite pairing traces]
\lean{FABL.PartitePairingTrace}
\label{FABL.PartitePairingTrace}
\leanok
\uses{FABL.BookPartiteVertex, FABL.greedyPartnerChoices, FABL.greedyPivot}
Let $r,k$ be natural numbers with $r>0$. For each natural number $s$ and each finite set
of vertices in a $2r$-partite vertex set, define the type of pairing traces recursively:
at $s=0$ it has one element; at $s+1$, a trace consists of the condition that the
current set is nonempty, a partner choice from the vertices in parts other than the
pivot's part, and a trace of length $s$ after removing the pivot and that partner.
\end{definition}

\begin{definition}[Finite pairing trace types]
\lean{FABL.partitePairingTraceFintype}
\label{FABL.partitePairingTraceFintype}
\leanok
\uses{FABL.BookPartiteVertex, FABL.PartitePairingTrace, FABL.greedyPartnerChoices, FABL.greedyPivot}
Let $r,k$ be natural numbers with $r>0$. For each natural number $s$ and each finite set
of vertices in a $2r$-partite vertex set, the pairing trace type of length $s$ is
equipped with a finite type structure.
\end{definition}

\begin{definition}[Successor trace decomposition]
\lean{FABL.partitePairingTraceSuccEquiv}
\label{FABL.partitePairingTraceSuccEquiv}
\leanok
\uses{FABL.BookPartiteVertex, FABL.PartitePairingTrace, FABL.greedyPartnerChoices, FABL.greedyPivot}
Let $r,k,s$ be natural numbers with $r>0$, and let a nonempty finite set of vertices be
given in a $2r$-partite vertex set. Pairing traces of length $s+1$ are equivalent to
pairs consisting of a partner choice from the vertices in parts other than the pivot's
part and a pairing trace of length $s$ on the set remaining after removing the pivot and
that partner.
\end{definition}

\begin{theorem}[Card partite pairing trace le]
\lean{FABL.card_partitePairingTrace_le}
\statementsource{Analysis-of-Boolean-Functions-by-Ryan-ODonnell}{
  pdf-ffbe9c7b3306-p277-b004,
  pdf-ffbe9c7b3306-p278-b001
}
\label{FABL.card_partitePairingTrace_le}
\leanok
\uses{FABL.BookPartiteVertex, FABL.PartitePairingTrace, FABL.card_greedyPartnerChoices_le, FABL.greedyPartnerChoices, FABL.greedyPivot, FABL.greedyPivot_mem, FABL.partitePairingTraceSuccEquiv}
\difficulty{5}
% dataset-pipeline: stub -- prose pending build_informal_statements.py
Statement pending regeneration.
\end{theorem}

\begin{definition}[Partner in a partite matching]
\lean{FABL.graphPartiteMatchingPartner}
\statementsource{Analysis-of-Boolean-Functions-by-Ryan-ODonnell}{pdf-ffbe9c7b3306-p277-b004}
\label{FABL.graphPartiteMatchingPartner}
\leanok
\uses{FABL.BookPartiteVertex, FABL.GraphPartitePerfectMatching}
Let $q$ and $k$ be natural numbers, and let a perfect matching be given in the complete
multipartite graph with $q$ parts of size $k$. The partner of a vertex is the vertex
paired with it by the matching.
\end{definition}

\begin{theorem}[Adjacency of the matching partner]
\lean{FABL.graphPartiteMatchingPartner_adj}
\label{FABL.graphPartiteMatchingPartner_adj}
\leanok
\uses{FABL.BookPartiteVertex, FABL.GraphPartitePerfectMatching, FABL.graphPartiteMatchingPartner}
\difficulty{1}
Let $q,k$ be nonnegative integers, and let $M$ be a perfect matching in the complete
$q$-partite graph with $k$ vertices in each part. For every vertex $v$, its matching
partner is adjacent to $v$ in the graph $M$.
\end{theorem}

\begin{theorem}[Uniqueness of the matching partner]
\lean{FABL.graphPartiteMatchingPartner_unique}
\label{FABL.graphPartiteMatchingPartner_unique}
\leanok
\uses{FABL.BookPartiteVertex, FABL.GraphPartitePerfectMatching, FABL.graphPartiteMatchingPartner}
\difficulty{1}
Let $q,k$ be nonnegative integers, and let $M$ be a perfect matching in the complete
$q$-partite graph with $k$ vertices in each part. If a vertex $w$ is adjacent to a
vertex $v$ in $M$, then $w$ is the matching partner of $v$.
\end{theorem}

\begin{theorem}[Involutive matching partner map]
\lean{FABL.graphPartiteMatchingPartner_involutive}
\label{FABL.graphPartiteMatchingPartner_involutive}
\leanok
\uses{FABL.GraphPartitePerfectMatching, FABL.graphPartiteMatchingPartner, FABL.graphPartiteMatchingPartner_adj, FABL.graphPartiteMatchingPartner_unique}
\difficulty{1}
Let $q,k$ be nonnegative integers, and let $M$ be a perfect matching in the complete
$q$-partite graph with $k$ vertices in each part. The map sending each vertex to its
matching partner is an involution: applying it twice returns the original vertex.
\end{theorem}

\begin{theorem}[Matching partners lie in different parts]
\lean{FABL.graphPartiteMatchingPartner_part_ne}
\label{FABL.graphPartiteMatchingPartner_part_ne}
\leanok
\uses{FABL.BookPartiteVertex, FABL.GraphPartitePerfectMatching, FABL.graphPartiteMatchingPartner, FABL.graphPartiteMatchingPartner_adj}
\difficulty{3}
Let $q,k$ be nonnegative integers, and let $M$ be a perfect matching in the complete
$q$-partite graph with $k$ vertices in each part. Every vertex and its matching partner
belong to different parts.
\end{theorem}

\begin{definition}[Closure under matching partners]
\lean{FABL.GraphMatchingClosedOn}
\statementsource{Analysis-of-Boolean-Functions-by-Ryan-ODonnell}{pdf-ffbe9c7b3306-p277-b004}
\label{FABL.GraphMatchingClosedOn}
\leanok
\uses{FABL.BookPartiteVertex, FABL.GraphPartitePerfectMatching, FABL.graphPartiteMatchingPartner}
Let $q$ and $k$ be natural numbers, let a perfect matching be given in the complete
multipartite graph with $q$ parts of size $k$, and let $R$ be a set of vertices. The set
$R$ is closed under matching partners if the partner of every vertex in $R$ also belongs
to $R$.
\end{definition}

\begin{theorem}[Closure of the full vertex set]
\lean{FABL.graphMatchingClosedOn_univ}
\label{FABL.graphMatchingClosedOn_univ}
\leanok
\uses{FABL.GraphMatchingClosedOn, FABL.GraphPartitePerfectMatching}
\difficulty{1}
In a complete multipartite graph with $q$ parts of $k$ vertices each and a perfect
matching, the full vertex set contains the matching partner of every vertex.
\end{theorem}

\begin{theorem}[Closure after removing a matched pair]
\lean{FABL.graphMatchingClosedOn_erase_pair}
\label{FABL.graphMatchingClosedOn_erase_pair}
\leanok
\uses{FABL.BookPartiteVertex, FABL.GraphMatchingClosedOn, FABL.GraphPartitePerfectMatching, FABL.graphPartiteMatchingPartner, FABL.graphPartiteMatchingPartner_involutive}
\difficulty{4}
Let a complete multipartite graph have vertex set consisting of $q$ parts with $k$
vertices each, and let it have a perfect matching. If a finite set of vertices contains
the matching partner of each of its vertices, then after removing a vertex in that set
and its matching partner, the remaining set still contains the matching partner of each
of its vertices.
\end{theorem}

\begin{definition}[Recursive graph matching trace]
\lean{FABL.graphMatchingTrace}
\statementsource{Analysis-of-Boolean-Functions-by-Ryan-ODonnell}{pdf-ffbe9c7b3306-p277-b004}
\label{FABL.graphMatchingTrace}
\leanok
\uses{FABL.BookPartiteVertex, FABL.GraphMatchingClosedOn, FABL.GraphPartitePerfectMatching, FABL.PartitePairingTrace, FABL.graphMatchingClosedOn_erase_pair, FABL.graphPartiteMatchingPartner, FABL.graphPartiteMatchingPartner_part_ne, FABL.greedyPartnerChoices, FABL.greedyPivot, FABL.greedyPivot_mem, FABL.partitePairingTraceSuccEquiv}
Let $r,k,s\in\mathbb{N}$ with $r>0$, and let $M$ be a perfect matching of the complete
$2r$-partite graph whose parts each have $k$ vertices. Suppose a set $R$ of vertices has
cardinality $2s$ and is closed under the matching partner map, meaning that the partner
of every vertex in $R$ also lies in $R$. Then one can form a pairing trace of length $s$
on $R$: at each step, choose a pivot in a largest part of the remaining set and a
partner from a different part, then continue after deleting both vertices.
\end{definition}

\begin{theorem}[Partners determined by the pairing trace]
\lean{FABL.graphMatchingTrace_partner_eq_on_remaining}
\label{FABL.graphMatchingTrace_partner_eq_on_remaining}
\leanok
\uses{FABL.BookPartiteVertex, FABL.GraphMatchingClosedOn, FABL.GraphPartitePerfectMatching, FABL.PartitePairingTrace, FABL.graphMatchingClosedOn_erase_pair, FABL.graphMatchingTrace, FABL.graphPartiteMatchingPartner, FABL.graphPartiteMatchingPartner_involutive, FABL.graphPartiteMatchingPartner_part_ne, FABL.greedyPartnerChoices, FABL.greedyPivot, FABL.greedyPivot_mem, FABL.partitePairingTraceSuccEquiv, FABL.proof_heq}
\difficulty{7}
Let $r,k\in\mathbb{N}$ with $r>0$, and let $R$ be a set of vertices in the complete
$2r$-partite graph with $|R|=2s$. Suppose two perfect matchings each send every vertex
of $R$ to its partner within $R$, and their greedy pairing traces on $R$ agree. Then the
two matchings assign the same partner to every vertex of $R$.
\end{theorem}

\begin{definition}[Full graph matching trace]
\lean{FABL.graphMatchingFullTrace}
\statementsource{Analysis-of-Boolean-Functions-by-Ryan-ODonnell}{pdf-ffbe9c7b3306-p277-b004}
\label{FABL.graphMatchingFullTrace}
\leanok
\uses{FABL.BookPartiteVertex, FABL.GraphPartitePerfectMatching, FABL.PartitePairingTrace, FABL.graphMatchingClosedOn_univ, FABL.graphMatchingTrace}
Let $r,k\in\mathbb{N}$ with $r>0$, and let $M$ be a perfect matching of the complete
$2r$-partite graph whose parts each have $k$ vertices. The entire vertex set admits a
pairing trace of length $rk$, in which each step chooses a pivot in a largest remaining
part and a partner from a different part, and deletes both vertices.
\end{definition}

\begin{theorem}[Injectivity of the full matching trace]
\lean{FABL.graphMatchingFullTrace_injective}
\label{FABL.graphMatchingFullTrace_injective}
\leanok
\uses{FABL.BookPartiteVertex, FABL.GraphPartitePerfectMatching, FABL.PartitePairingTrace, FABL.graphMatchingClosedOn_univ, FABL.graphMatchingFullTrace, FABL.graphMatchingTrace_partner_eq_on_remaining, FABL.graphPartiteMatchingPartner, FABL.graphPartiteMatchingPartner_adj, FABL.graphPartiteMatchingPartner_unique}
\difficulty{4}
Let $r,k\in\mathbb{N}$ with $r>0$, and consider the complete $2r$-partite graph with $k$
vertices in each part. The map sending a perfect matching to its full greedy pairing
trace, starting from all vertices and making $rk$ pairings, is injective.
\end{theorem}

\begin{theorem}[Card partite perfect matching le factorial bound strong]
\lean{FABL.card_partitePerfectMatching_le_factorial_bound_strong}
\statementsource{Analysis-of-Boolean-Functions-by-Ryan-ODonnell}{
  pdf-ffbe9c7b3306-p277-b004,
  pdf-ffbe9c7b3306-p278-b001
}
\label{FABL.card_partitePerfectMatching_le_factorial_bound_strong}
\leanok
\uses{FABL.BookPartiteVertex, FABL.GraphPartitePerfectMatching, FABL.PartitePairingTrace, FABL.PartitePerfectMatching, FABL.card_graphPartitePerfectMatching_eq_partitePerfectMatching, FABL.card_partitePairingTrace_le, FABL.graphMatchingFullTrace, FABL.graphMatchingFullTrace_injective}
\difficulty{3}
% dataset-pipeline: stub -- prose pending build_informal_statements.py
Statement pending regeneration.
\end{theorem}

\begin{theorem}[Card partite perfect matching le factorial bound printed]
\lean{FABL.card_partitePerfectMatching_le_factorial_bound_printed}
\statementsource{Analysis-of-Boolean-Functions-by-Ryan-ODonnell}{
  pdf-ffbe9c7b3306-p277-b004,
  pdf-ffbe9c7b3306-p278-b001
}
\label{FABL.card_partitePerfectMatching_le_factorial_bound_printed}
\leanok
\uses{FABL.PartitePerfectMatching, FABL.card_partitePerfectMatching_le_factorial_bound_strong}
\difficulty{3}
% dataset-pipeline: stub -- prose pending build_informal_statements.py
Statement pending regeneration.
\end{theorem}

\begin{theorem}[Factorial matching count identity]
\lean{FABL.matching_count_root_div_sqrt_factorial}
\label{FABL.matching_count_root_div_sqrt_factorial}
\leanok
\uses{FABL.matchingImprovementConstant}
\difficulty{6}
For positive integers $r$ and natural numbers $k$, the $1/(2r)$-th power of
$\bigl((2r-1)/r\bigr)^{rk}(rk)!$, divided by $\sqrt{k!}$, equals the matching
improvement constant times $(\sqrt{2r-1})^k$.
\end{theorem}

\begin{theorem}[Improved homogeneous moment bound]
\lean{FABL.homogeneous_uniformLpNorm_le_matchingImprovement}
\proofsource{Analysis-of-Boolean-Functions-by-Ryan-ODonnell}{
  pdf-ffbe9c7b3306-p277-b004,
  pdf-ffbe9c7b3306-p278-b001
}
\label{FABL.homogeneous_uniformLpNorm_le_matchingImprovement}
\leanok
\uses{FABL.IsFourierHomogeneous, FABL.PartitePerfectMatching, FABL.card_partitePerfectMatching_le_factorial_bound_strong, FABL.homogeneous_uniformLpNorm_le_partiteMatching, FABL.matchingImprovementConstant, FABL.matching_count_root_div_sqrt_factorial, FABL.uniformLpNorm, FABL.uniformLpNorm_nonneg}
\difficulty{4}
Let $n,r,k$ be positive integers, and let $f$ be a real-valued function on the sign cube
$\{-1,1\}^n$ whose nonzero Fourier coefficients are all supported on sets of size $k$.
Then its uniform $2r$ norm is at most
\[C_{r,k}\,\bigl(\sqrt{2r-1}\bigr)^k\,\|f\|_2,\]
where
\[C_{r,k}=\left(\frac{(rk)!}{(k!)^r r^{rk}}\right)^{1/(2r)}.\]
\end{theorem}

\begin{theorem}[Homogeneous uniform lp norm le sqrt two mul sub one]
\lean{FABL.homogeneous_uniformLpNorm_le_sqrt_two_mul_sub_one}
\statementsource{Analysis-of-Boolean-Functions-by-Ryan-ODonnell}{
  pdf-ffbe9c7b3306-p277-b004,
  pdf-ffbe9c7b3306-p278-b001
}
\label{FABL.homogeneous_uniformLpNorm_le_sqrt_two_mul_sub_one}
\leanok
\uses{FABL.IsFourierHomogeneous, FABL.homogeneous_uniformLpNorm_le_matchingImprovement, FABL.matchingImprovementConstant, FABL.matchingImprovementConstant_le_one, FABL.uniformLpNorm, FABL.uniformLpNorm_nonneg}
\difficulty{4}
% dataset-pipeline: stub -- prose pending build_informal_statements.py
Statement pending regeneration.
\end{theorem}

\begin{theorem}[Homogeneous moment norm bound]
\lean{FABL.homogeneous_uniformLpNorm_le_doubleFactorial}
\proofsource{Analysis-of-Boolean-Functions-by-Ryan-ODonnell}{
  pdf-ffbe9c7b3306-p277-b004,
  pdf-ffbe9c7b3306-p278-b001
}
\label{FABL.homogeneous_uniformLpNorm_le_doubleFactorial}
\leanok
\uses{FABL.IsFourierHomogeneous, FABL.PartitePerfectMatching, FABL.homogeneous_uniformLpNorm_le_partiteMatching, FABL.natCard_partitePerfectMatching_le_doubleFactorial, FABL.uniformLpNorm, FABL.uniformLpNorm_nonneg}
\difficulty{4}
Let $n,r,k$ be positive integers, and let $f$ be a real-valued function on the sign cube
$\{-1,1\}^n$ whose nonzero Fourier coefficients are all supported on sets of size $k$.
Then the uniform $2r$ norm of $f$ is at most $\bigl((2rk-1)!!\bigr)^{1/(2r)}$ and at
most $\sqrt{k!}$ times the uniform second norm of $f$.
\end{theorem}

\begin{theorem}[Is fourier homogeneous fourier degree le]
\lean{FABL.isFourierHomogeneous_fourierDegree_le}
\statementsource{Analysis-of-Boolean-Functions-by-Ryan-ODonnell}{pdf-ffbe9c7b3306-p036-b002}
\label{FABL.isFourierHomogeneous_fourierDegree_le}
\leanok
\uses{FABL.IsFourierHomogeneous, FABL.fourierDegree, FABL.fourierDegree_le_iff}
\difficulty{3}
% dataset-pipeline: stub -- prose pending build_informal_statements.py
Statement pending regeneration.
\end{theorem}

\begin{theorem}[Homogeneous uniform lp norm le sqrt q]
\lean{FABL.homogeneous_uniformLpNorm_le_sqrt_q}
\statementsource{Analysis-of-Boolean-Functions-by-Ryan-ODonnell}{
  pdf-ffbe9c7b3306-p247-b006,
  pdf-ffbe9c7b3306-p277-b004,
  pdf-ffbe9c7b3306-p278-b001
}
\label{FABL.homogeneous_uniformLpNorm_le_sqrt_q}
\leanok
\uses{FABL.IsFourierHomogeneous, FABL.homogeneous_uniformLpNorm_le_sqrt_two_mul, FABL.isFourierHomogeneous_fourierDegree_le, FABL.uniformLpNorm}
\difficulty{1}
% dataset-pipeline: stub -- prose pending build_informal_statements.py
Statement pending regeneration.
\end{theorem}

\begin{theorem}[Strict factorial bound]
\lean{FABL.factorial_bound_strong_lt_printed}
\label{FABL.factorial_bound_strong_lt_printed}
\leanok
\difficulty{4}
Let $r$ and $k$ be natural numbers with $r>0$ and $1<rk$. Then
$\bigl((2r-1)/r\bigr)^{rk}(rk)!<\bigl((2r-1)/r\bigr)^{rk}((rk)!)^2$.
\end{theorem}

\begin{definition}[Six blocks yield a parity tuple]
\lean{FABL.SixDisjointBlockChoice.toParityTuple}
\statementsource{Analysis-of-Boolean-Functions-by-Ryan-ODonnell}{pdf-ffbe9c7b3306-p277-b004}
\label{FABL.SixDisjointBlockChoice.toParityTuple}
\leanok
\uses{FABL.LevelKParityTuple, FABL.SixDisjointBlockChoice, FABL.SixDisjointBlockChoice.incidentSupport, FABL.SixDisjointBlockChoice.incidentSupport_card, FABL.SixDisjointBlockChoice.tupleSymmDiff_incidentSupport}
For natural numbers $n$ and $d$, a choice of six ordered, pairwise disjoint $d$-element
subsets of $\{0,\ldots,n-1\}$ determines a tuple of four subsets of $\{0,\ldots,n-1\}$.
Each of the four sets has size $3d$, and every element belongs to an even number of the
four sets.
\end{definition}

\begin{theorem}[Injectivity of the parity tuple map]
\lean{FABL.SixDisjointBlockChoice.toParityTuple_injective}
\label{FABL.SixDisjointBlockChoice.toParityTuple_injective}
\leanok
\uses{FABL.DisjointBlockChoices.blocks_length, FABL.DisjointBlockChoices.ext_of_blocks_eq, FABL.LevelKParityTuple, FABL.SixDisjointBlockChoice, FABL.SixDisjointBlockChoice.block, FABL.SixDisjointBlockChoice.block_eq_incidentSupport_inter, FABL.SixDisjointBlockChoice.incidentSupport, FABL.SixDisjointBlockChoice.toParityTuple, FABL.k4EdgeEnds}
\difficulty{4}
For fixed natural numbers $n$ and $d$, the map sending a choice of six pairwise disjoint
$d$-element subsets of an $n$-element set to the four incident supports indexed by the
vertices of the complete graph on four vertices is injective.
\end{theorem}

\begin{theorem}[Expect level k polynomial pow eq nat card parity tuple]
\lean{FABL.expect_levelKPolynomial_pow_eq_natCard_parityTuple}
\statementsource{Analysis-of-Boolean-Functions-by-Ryan-ODonnell}{pdf-ffbe9c7b3306-p277-b004}
\label{FABL.expect_levelKPolynomial_pow_eq_natCard_parityTuple}
\leanok
\uses{FABL.LevelKParityTuple, FABL.expect_pow_eq_sum_fourierCoeff_of_tupleSymmDiff_empty, FABL.fourierCoeff_levelKPolynomial, FABL.levelKPolynomial, FABL.tupleSymmDiff}
\difficulty{5}
% dataset-pipeline: stub -- prose pending build_informal_statements.py
Statement pending regeneration.
\end{theorem}

\begin{theorem}[Fourth moment bounds the number of encoded tuples]
\lean{FABL.natCard_le_fourthMoment_levelKPolynomial_of_injective}
\label{FABL.natCard_le_fourthMoment_levelKPolynomial_of_injective}
\leanok
\uses{FABL.LevelKParityTuple, FABL.expect_levelKPolynomial_pow_eq_natCard_parityTuple, FABL.levelKPolynomial}
\difficulty{1}
Let $W$ be a finite type, and let $n,k\in\mathbb N$. If there is an injective map from
$W$ into the set of ordered quadruples of subsets of $[n]$, each of size $k$, whose
symmetric difference is empty, then $|W|$ is at most the uniform average over
$x\in\{−1,1\}^n$ of the fourth power of the level-$k$ polynomial evaluated at $x$.
\end{theorem}

\begin{theorem}[Six-block multinomial moment bound]
\lean{FABL.sixBlockMultinomial_le_fourthMoment_levelKPolynomial}
\label{FABL.sixBlockMultinomial_le_fourthMoment_levelKPolynomial}
\leanok
\uses{FABL.SixDisjointBlockChoice, FABL.SixDisjointBlockChoice.toParityTuple, FABL.SixDisjointBlockChoice.toParityTuple_injective, FABL.levelKPolynomial, FABL.natCard_le_fourthMoment_levelKPolynomial_of_injective, FABL.natCard_sixDisjointBlockChoice_eq_multinomial, FABL.sixBlockMultinomial}
\difficulty{3}
Let $n$ and $d$ be nonnegative integers with $6d\\le n$. The multinomial coefficient
with six parts of size $d$ and a seventh part of size $n-6d$ is at most the fourth
moment, under the uniform probability distribution on the $n$-dimensional sign cube, of
the sum of all degree-$3d$ monomials.
\end{theorem}

\begin{theorem}[Fourth moment lower bound]
\lean{FABL.exercise9_3_fourthMoment_lower_bound}
\label{FABL.exercise9_3_fourthMoment_lower_bound}
\leanok
\uses{FABL.expect_levelKPolynomial_sq, FABL.levelKPolynomial, FABL.sixBlockMultinomial, FABL.sixBlockMultinomial_le_fourthMoment_levelKPolynomial}
\difficulty{4}
Let $n,d$ be nonnegative integers with $6d\leq n$, and let $F$ be the level-$3d$
polynomial on the discrete cube of dimension $n$. If $M(n,d)$ denotes the multinomial
coefficient with six entries equal to $d$ and a seventh equal to $n-6d$, then
\[
\frac{M(n,d)}{\binom{n}{3d}^{2}}\bigl(\mathbb{E}F^2\bigr)^2\leq \mathbb{E}F^4,
\]
where expectations are uniform over the cube.
\end{theorem}

\begin{theorem}[Fourth moment bound for divisible degrees]
\lean{FABL.exercise9_3_fourthMoment_lower_bound_of_dvd}
\label{FABL.exercise9_3_fourthMoment_lower_bound_of_dvd}
\leanok
\uses{FABL.exercise9_3_fourthMoment_lower_bound, FABL.levelKPolynomial, FABL.sixBlockMultinomial}
\difficulty{4}
Let $n,k$ be nonnegative integers such that $k>0$, $3$ divides $k$, and $2k\leq n$. Let
$F$ be the level-$k$ polynomial on the discrete cube of dimension $n$, and let
$M(n,k/3)$ denote the multinomial coefficient with six entries equal to $k/3$ and a
seventh equal to $n-2k$. Then
\[
\frac{M(n,k/3)}{\binom{n}{k}^{2}}\bigl(\mathbb{E}F^2\bigr)^2\leq \mathbb{E}F^4,
\]
where expectations are uniform over the cube.
\end{theorem}
